\documentclass[11pt,oneside]{article}

% Standalone ProveIt research article.  The source intentionally has no
% external inputs so that the TeX/PDF pair is directly reproducible after
% extraction.
\usepackage[a4paper,margin=26mm,headheight=15pt]{geometry}
\usepackage[T1]{fontenc}
\usepackage[utf8]{inputenc}
\IfFileExists{libertinus.sty}{\usepackage{libertinus}}{\usepackage{lmodern}}
\usepackage{microtype}
\usepackage{amsmath,amssymb,amsthm,mathtools,mathrsfs,bm}
\usepackage{booktabs,longtable,tabularx,array,multirow}
\usepackage{graphicx}
\usepackage{xcolor}
\usepackage{enumitem}
\usepackage{fancyhdr}
\usepackage{xurl}
\usepackage{listings}
\usepackage[most]{tcolorbox}
\usepackage{hyperref}
\usepackage[nameinlink,noabbrev,capitalise]{cleveref}

\definecolor{linkblue}{RGB}{24,72,124}
\definecolor{deepgreen}{RGB}{23,102,69}
\definecolor{deepred}{RGB}{140,42,42}
\definecolor{deepgold}{RGB}{122,91,19}
\definecolor{shadegray}{RGB}{246,247,249}
\definecolor{rulegray}{RGB}{195,201,208}
\definecolor{palered}{RGB}{251,236,236}
\definecolor{palegold}{RGB}{251,245,230}
\definecolor{palegreen}{RGB}{233,245,243}
\definecolor{paleblue}{RGB}{236,243,251}

\hypersetup{
  hypertexnames=false,
  colorlinks=true,
  linkcolor=linkblue,
  citecolor=deepgreen,
  urlcolor=linkblue,
  pdftitle={Asymptotic Inversion of the Generalized Hyperfactorial K-Function},
  pdfauthor={Vladimir Reshetnikov},
  pdfsubject={Inverse asymptotics for the generalized hyperfactorial K-function and a general theory of power-logarithmic transseries inversion},
  pdfkeywords={K-function, generalized hyperfactorial, hyperfactorial, Lambert W, r-Lambert W, inverse asymptotics, transseries, Bernoulli polynomials, Barnes G-function, Hurwitz zeta}
}

\pagestyle{fancy}
\fancyhf{}
\fancyhead[L]{\small Inverse transseries of the \(K\)-function}
\fancyhead[R]{\small\nouppercase{\leftmark}}
\fancyfoot[C]{\thepage}
\renewcommand{\headrulewidth}{0.4pt}

\setcounter{tocdepth}{2}
\setcounter{secnumdepth}{3}
\setlist[itemize]{leftmargin=1.5em,itemsep=0.25ex,topsep=0.6ex}
\setlist[enumerate]{leftmargin=1.9em,itemsep=0.35ex,topsep=0.6ex}
\allowdisplaybreaks[2]

\theoremstyle{plain}
\newtheorem{theorem}{Theorem}[section]
\newtheorem{proposition}[theorem]{Proposition}
\newtheorem{lemma}[theorem]{Lemma}
\newtheorem{corollary}[theorem]{Corollary}
\theoremstyle{definition}
\newtheorem{definition}[theorem]{Definition}
\newtheorem{algorithm}[theorem]{Algorithm}
\newtheorem{example}[theorem]{Example}
\newtheorem{conjecture}[theorem]{Conjecture}
\theoremstyle{remark}
\newtheorem{remark}[theorem]{Remark}

\crefname{algorithm}{algorithm}{algorithms}
\Crefname{algorithm}{Algorithm}{Algorithms}

\newtcolorbox{mainbox}[1][]{%
  enhanced,breakable,colback=paleblue,colframe=linkblue,
  boxrule=0.7pt,arc=1.2mm,left=2mm,right=2mm,top=1.4mm,bottom=1.4mm,
  title={Main result},fonttitle=\bfseries,#1}
\newtcolorbox{methodbox}[1][]{%
  enhanced,breakable,colback=palegreen,colframe=deepgreen,
  boxrule=0.6pt,arc=1.2mm,left=2mm,right=2mm,top=1.2mm,bottom=1.2mm,
  title={Reusable method},fonttitle=\bfseries,#1}
\newtcolorbox{warningbox}[1][]{%
  enhanced,breakable,colback=palered,colframe=deepred,
  boxrule=0.6pt,arc=1.2mm,left=2mm,right=2mm,top=1.2mm,bottom=1.2mm,
  title={Scale warning},fonttitle=\bfseries,#1}
\newtcolorbox{summarybox}[1][]{%
  enhanced,breakable,colback=palegold,colframe=deepgold,
  boxrule=0.6pt,arc=1.2mm,left=2mm,right=2mm,top=1.2mm,bottom=1.2mm,
  title={Summary},fonttitle=\bfseries,#1}
\newtcolorbox{statusbox}[1][]{%
  enhanced,breakable,colback=shadegray,colframe=rulegray,
  boxrule=0.6pt,arc=1.2mm,left=2mm,right=2mm,top=1.2mm,bottom=1.2mm,
  title={Research status},fonttitle=\bfseries,#1}

\lstdefinestyle{pythonstyle}{%
  language=Python,
  basicstyle=\ttfamily\small,
  columns=fullflexible,
  frame=single,
  rulecolor=\color{rulegray},
  backgroundcolor=\color{shadegray},
  showstringspaces=false,
  keepspaces=true,
  breaklines=true,
  xleftmargin=0.4em,xrightmargin=0.4em,
  aboveskip=0.8em,belowskip=0.8em
}
\lstset{style=pythonstyle}

\newcommand{\Kfun}{\operatorname{K}}
\newcommand{\Gfun}{\operatorname{G}}
\newcommand{\W}{\operatorname{W}}
\newcommand{\rW}{\mathcal{W}}
\newcommand{\e}{\mathrm{e}}
\newcommand{\dd}{\mathop{}\!\mathrm{d}}
\newcommand{\R}{\mathbb{R}}
\newcommand{\Q}{\mathbb{Q}}
\newcommand{\C}{\mathbb{C}}
\newcommand{\N}{\mathbb{N}}
\newcommand{\B}{\mathrm{B}}
\newcommand{\BigO}{\mathcal{O}}
\newcommand{\smallo}{\mathop{o}}
\newcommand{\coeff}[2]{\left[#1^{#2}\right]}
\newcommand{\inv}{^{-1}}
\newcommand{\calD}{\mathcal{D}}
\newcommand{\calE}{\mathcal{E}}
\newcommand{\calA}{\mathcal{A}}
\newcommand{\calS}{\mathcal{S}}
\newcommand{\calT}{\mathcal{T}}
\newcommand{\angles}[1]{\left\langle #1\right\rangle}
\DeclareMathOperator{\Res}{Res}
\DeclareMathOperator{\Trunc}{Trunc}
\DeclareMathOperator{\val}{val}

\begin{document}

\title{\textbf{Asymptotic Inversion of the Generalized Hyperfactorial \(K\)-Function}\\[0.4em]
  \Large Lambert and \(r\)-Lambert Anchors, Centered Bernoulli Structure,\\
  and a General Calculus for Power--Logarithmic Transseries}
\author{Vladimir Reshetnikov\\
  \small ProveIt research note}
\date{3 September 2026}
\maketitle

\begin{abstract}
The generalized hyperfactorial, or \(K\)-function, is characterized by
\(\Kfun(z+1)=z^z\Kfun(z)\), \(\Kfun(1)=1\), and admits the exact representations
\[
 \log\Kfun(z)=\zeta'(-1,z)-\zeta'(-1)
 =(z-1)\log\Gamma(z)-\log\Gfun(z).
\]
This article derives an all-orders asymptotic transseries for the large positive
inverse branch \(\Kfun_+^{-1}(y)\) as \(y\to+\infty\), and then develops a
general inversion theory for functions whose logarithms have dominant shape
\(a x^p(\log x-\beta)\).

The decisive normalization is the midpoint shift \(x=q+\tfrac12\).  It removes
the otherwise present \(q\log q\) term and turns the entire algebraic tail into
an even series:
\[
 \log\Kfun\!\left(q+\tfrac12\right)
 \sim \frac{q^2}{2}\log q-\frac{q^2}{4}
      -\frac1{24}\log q+\kappa
      +\sum_{m\ge1}\frac{c_m}{q^{2m}},
 \qquad \kappa=\log A-\frac18,
\]
where \(A\) is the Glaisher--Kinkelin constant and
\[
 c_m=-\frac{B_{2m+2}(1/2)}{(2m)(2m+1)(2m+2)}.
\]
For \(Y=\log y\), put
\[
 w=\W_0\!\left(\frac{4Y}{\e}\right),\qquad
 Q=2\sqrt{\frac{Y}{w}}=\exp\!\left(\frac{w+1}{2}\right),
 \qquad L=\log Q.
\]
Then
\[
 \Kfun_+^{-1}(y)
 \sim \frac12+Q+\frac{a_1(L)}{Q}
 +\frac{a_2(L)}{Q^3}+\frac{a_3(L)}{Q^5}+\cdots,
\]
with explicit rational coefficient functions, a triangular coefficient
recurrence, and an operator formula valid to all orders.  A refined anchor
absorbs the logarithmic defect into an exactly invertible core.  It is expressed
through the \(r\)-Lambert function with \(r=-1/(12\e)\), and moves the first
correction from order \(Q^{-1}\) to order \(Q^{-3}/\log Q\).

For the general class
\[
 F(x)\sim a x^p(\log x-\beta)
       +\sum_{\lambda\in\Lambda}x^{p-\lambda}r_\lambda(U),
 \qquad U=p(\log x-\beta),
\]
we prove an anchored Lagrange--B\"urmann formula, a support-transfer theorem,
and a closed all-orders expression for every inverse coefficient.  A forward
defect of power gap \(\lambda\) produces an inverse defect beginning at power
\(x^{1-\lambda}\), while mixed interactions are indexed by the additive
semigroup generated by the forward gaps.  The theory explains, rather than
merely observes, the odd inverse-power structure of the \(K\)-function.
Numerical tests against high-precision Barnes-\(G\) evaluation confirm the
predicted orders and show that the refined anchor is already highly accurate at
moderate arguments.
\end{abstract}

\noindent\textbf{Keywords.}
Generalized hyperfactorial; \(K\)-function; inverse asymptotics; transseries;
Lambert \(W\); \(r\)-Lambert function; Hurwitz zeta; Barnes \(G\)-function;
Bernoulli polynomials; Lagrange--B\"urmann inversion.

\begin{statusbox}
Exact functional identities and the real-branch derivative formulas below are
classical.  The midpoint normal form follows rigorously from the standard
Hurwitz-zeta expansion.  The inverse coefficient formulas and recurrences are
formal consequences of Lagrange--B\"urmann inversion and become Poincar\'e
asymptotic expansions on the positive real ray by the remainder estimates in
\cref{sec:analytic-validity}.  Statements about exponentially improved complex
expansions are explicitly separated from the algebraic theory.
\end{statusbox}

{\small\tableofcontents}
\clearpage

\section{The problem and the answer at a glance}
\label{sec:overview}

The discrete hyperfactorial
\[
 H_n=\prod_{k=1}^{n}k^k
\]
has logarithm of order \(n^2\log n\).  The analytic continuation relevant here
is the \(K\)-function, normalized so that
\begin{equation}
 \Kfun(n+1)=H_n,\qquad n\in\N.
 \label{eq:integer-hyperfactorial}
\end{equation}
The inverse problem is not to invert \(\Kfun\) by an elementary operation, but
to understand the complete hierarchy of scales in the solution of
\begin{equation}
 \Kfun(x)=y,
 \qquad y\to+\infty.
 \label{eq:inverse-problem}
\end{equation}
Taking logarithms converts an extremely fast-growing function into one whose
dominant part is a quadratic logarithmic monomial.  Set
\begin{equation}
 Y=\log y.
 \label{eq:Ydef}
\end{equation}
The rough balance \(Y\sim \tfrac12 x^2\log x\) already suggests
\(x\asymp\sqrt{Y/\log Y}\).  However, a complete inverse expansion contains
three distinct kinds of information:
\begin{enumerate}
\item the nested logarithms coming from the ordinary Lambert function;
\item a constant centering correction \(1/2\), which is beyond every algebraic
      order in \(1/\log Y\) relative to the leading term;
\item an infinite hierarchy of algebraically small inverse-power sectors.
\end{enumerate}
A single ordinary power series cannot encode all three layers without losing
their ordering.

\begin{mainbox}
Let
\begin{equation}
 \kappa=\log A-\frac18,
 \qquad
 w=\W_0\!\left(\frac{4Y}{\e}\right),
 \qquad
 Q=\exp\!\left(\frac{w+1}{2}\right)
   =2\sqrt{\frac{Y}{w}},
 \qquad L=\log Q.
 \label{eq:main-anchor-data}
\end{equation}
Define the slow defect
\begin{equation}
 \Delta=\frac{L}{24}-\kappa.
 \label{eq:Deltadef}
\end{equation}
Then the large positive inverse branch has the Poincar\'e transseries
\begin{equation}
 \boxed{
 \Kfun_+^{-1}(y)
 \sim \frac12+Q+\sum_{n\ge1}\frac{a_n(L)}{Q^{2n-1}}.}
 \label{eq:main-inverse-series}
\end{equation}
The first three coefficient functions are
\begin{align}
 a_1(L)={}&\frac{\Delta}{L},
 \label{eq:a1}\displaybreak[0]\\
 a_2(L)={}&\frac{7}{5760L}+\frac{\Delta}{24L^2}
 -\frac{(L+1)\Delta^2}{2L^3},
 \label{eq:a2}\displaybreak[0]\\
 a_3(L)={}&
 \frac{(3L^2+5L+3)\Delta^3}{6L^5}
 -\frac{(4L+3)\Delta^2}{48L^4}
 +\frac{(1-7L)\Delta}{1920L^3}
 +\frac{49-186L}{967680L^2}.
 \label{eq:a3}
\end{align}
For every fixed \(N\ge0\),
\begin{equation}
 \Kfun_+^{-1}(y)
 =\frac12+Q+\sum_{n=1}^{N}\frac{a_n(L)}{Q^{2n-1}}
 +\BigO\!\left(Q^{-2N-1}\right).
 \label{eq:main-remainder}
\end{equation}
All coefficient functions are generated recursively by
\cref{thm:triangular-recurrence}, or directly by the support formula in
\cref{thm:support-transfer}.
\end{mainbox}

There is a second, sharper organization.  Let \(\rW_r\) denote the inverse of
\(s\mapsto s\e^s+rs\), and put
\begin{equation}
 s=\rW_{-1/(12\e)}\!\left(\frac{48(Y-\kappa)+1}{12\e}\right),
 \qquad Q_*=\exp\!\left(\frac{s+1}{2}\right),
 \qquad L_*=\log Q_*.
 \label{eq:refined-anchor-data}
\end{equation}
Then
\begin{equation}
 \boxed{
 \Kfun_+^{-1}(y)
 \sim \frac12+Q_*+\sum_{n\ge1}\frac{b_n(L_*)}{Q_*^{2n+1}},}
 \label{eq:refined-inverse-series}
\end{equation}
where
\begin{align}
 b_1(L)&=\frac{7}{5760L},
 \label{eq:b1}\displaybreak[0]\\
 b_2(L)&=-\frac{186L-49}{967680L^2},
 \label{eq:b2}\displaybreak[0]\\
 b_3(L)&=\frac{45720L^2-5435L+637}{464486400L^3}.
 \label{eq:b3}
\end{align}
The refined anchor incorporates the full defect of power gap two; consequently
its first unabsorbed forward term has gap four, and the inverse begins at
\(Q_*^{-3}\), not \(Q_*^{-1}\).

\begin{warningbox}
Expanding only in inverse powers of \(\log Y\) gives a correct but incomplete
view.  The term \(1/2\) and every inverse-power sector are beyond all orders of
that coarse logarithmic scale.  The phrase ``the asymptotic expansion'' is
therefore ambiguous unless the ordered transmonomial scale is stated.
\end{warningbox}

\section{The \texorpdfstring{\(K\)-function}{K-function} and its large real branch}
\label{sec:K-background}

\subsection{Equivalent definitions}

For \(\Re z>0\), one convenient definition is~\cite{Adamchik1998,DLMF25}
\begin{equation}
 \log\Kfun(z)=\zeta'(-1,z)-\zeta'(-1),
 \label{eq:K-Hurwitz}
\end{equation}
where \(\zeta(s,z)\) is the Hurwitz zeta function and the prime denotes the
derivative with respect to \(s\).  Equivalent formulas are
\begin{equation}
 \Kfun(z)=(2\pi)^{-(z-1)/2}
 \exp\!\left\{\binom{z}{2}+
       \int_0^{z-1}\log\Gamma(t+1)\,\dd t\right\},
 \label{eq:K-integral}
\end{equation}
and, using the Barnes \(G\)-function~\cite{DLMF5,Vardi1988},
\begin{equation}
 \Kfun(z)\Gfun(z)=\Gamma(z)^{z-1},
 \label{eq:K-Barnes}
\end{equation}
where \(\Gfun\) is the Barnes \(G\)-function normalized by
\(\Gfun(z+1)=\Gamma(z)\Gfun(z)\) and \(\Gfun(1)=1\).

The Hurwitz shift relation immediately gives
\begin{equation}
 \log\Kfun(z+1)-\log\Kfun(z)=z\log z,
 \label{eq:K-log-recurrence}
\end{equation}
so that
\begin{equation}
 \Kfun(z+1)=z^z\Kfun(z),\qquad \Kfun(1)=1.
 \label{eq:K-recurrence}
\end{equation}
Iterating \eqref{eq:K-recurrence} proves \eqref{eq:integer-hyperfactorial}.

\subsection{An exact logarithmic derivative}

The derivative of \eqref{eq:K-Hurwitz} is especially useful for inversion.
The identities
\[
 \partial_z\zeta(s,z)=-s\zeta(s+1,z),\qquad
 \zeta(0,z)=\frac12-z,
 \qquad
 \zeta'(0,z)=\log\Gamma(z)-\frac12\log(2\pi)
\]
yield the following formula.

\begin{proposition}[Exact derivative]
\label{prop:exact-derivative}
For \(x>0\),
\begin{align}
 \frac{\dd}{\dd x}\log\Kfun(x)
 & =\log\Gamma(x)+x-\frac12-\frac12\log(2\pi),
 \label{eq:dlogK}\displaybreak[0]\\
 \frac{\dd^2}{\dd x^2}\log\Kfun(x)
 & =\psi(x)+1,
 \label{eq:d2logK}
\end{align}
where \(\psi=\Gamma'/\Gamma\).
\end{proposition}

\begin{proof}
Differentiate \(\partial_z\zeta(s,z)=-s\zeta(s+1,z)\) with respect to \(s\):
\[
 \partial_z\zeta'(s,z)=-\zeta(s+1,z)-s\zeta'(s+1,z).
\]
At \(s=-1\), this becomes
\[
 \partial_z\zeta'(-1,z)=-\zeta(0,z)+\zeta'(0,z),
\]
and substitution of the two special values gives \eqref{eq:dlogK}.
A second differentiation gives \eqref{eq:d2logK}.
\end{proof}

Since \(\psi(x)=\log x+O(x^{-1})\), the function \(\log\Kfun(x)\) is strictly
increasing for all sufficiently large \(x\), and tends to \(+\infty\).  We use
\(\Kfun_+^{-1}\) for the inverse branch selected by \(x\to+\infty\) as
\(y\to+\infty\).

For orientation, these are the only two positive stationary points.
Indeed, \(\psi'(x)>0\), so \(\psi(x)+1\) crosses zero exactly once and
\((\log\Kfun)'\) has a unique minimum.  Moreover,
\((\log\Kfun)'(x)\to+\infty\) at both endpoints \(0^+\) and
\(+\infty\), while \((\log\Kfun)'(1)<0\).  The two zeros are
numerically
\begin{align*}
 x_{\max}&=0.2909569866991807959\ldots,
 &\Kfun(x_{\max})&=1.2986723069317708434\ldots,\\
 x_{\min}&=1.5376888637364865107\ldots,
 &\Kfun(x_{\min})&=0.8797868430593994168\ldots.
\end{align*}
Thus the notation \(\Kfun_+^{-1}\) is not cosmetic: below the upper extremal
value there are additional real local branches.  None survives in the limit
\(y\to+\infty\).

\section{Forward asymptotics and the midpoint normal form}
\label{sec:forward}

\subsection{The unshifted expansion}

Let \(A\) be the Glaisher--Kinkelin constant, characterized by
\begin{equation}
 \log A=\frac1{12}-\zeta'(-1).
 \label{eq:glaisher}
\end{equation}
The standard large-parameter expansion of the Hurwitz-zeta derivative~\cite{DLMF25,Nemes2017} gives,
for every fixed \(N\ge0\),
\begin{align}
 \log\Kfun(x)
 ={}&\left(\frac{x^2}{2}-\frac{x}{2}+\frac1{12}\right)\log x
     -\frac{x^2}{4}+\log A
 \notag\\
 &-\sum_{m=1}^{N}
   \frac{B_{2m+2}}{(2m)(2m+1)(2m+2)x^{2m}}
 +\BigO\!\left(x^{-2N-2}\right),
 \label{eq:unshifted-forward}
\end{align}
initially on the positive real ray, and uniformly in each closed subsector
\(|\arg x|\le\pi-\delta\) after a branch of \(\log x\) is fixed.  Here
\(B_n=B_n(0)\) are Bernoulli numbers.  The first terms are
\begin{equation}
 \log\Kfun(x)
 \sim \left(\frac{x^2}{2}-\frac{x}{2}+\frac1{12}\right)\log x
 -\frac{x^2}{4}+\log A
 +\frac1{720x^2}-\frac1{5040x^4}
 +\frac1{10080x^6}-\cdots.
 \label{eq:unshifted-first}
\end{equation}

The asymptotic form \eqref{eq:unshifted-forward} is already sufficient for a
leading inversion.  It is not, however, in the best coordinate.  The term
\(-\tfrac12x\log x\) creates both even and odd inverse powers after reversion.
A translation removes it and reveals a hidden Bernoulli symmetry.

\subsection{A shifted expansion valid for arbitrary center}

The Bernoulli polynomials are defined by
\begin{equation}
 \frac{te^{\alpha t}}{e^t-1}
 =\sum_{n=0}^{\infty}B_n(\alpha)\frac{t^n}{n!}.
 \label{eq:Bernoulli-polynomials}
\end{equation}
They satisfy
\begin{equation}
 B_n(\alpha)=\sum_{j=0}^{n}\binom{n}{j}B_j\alpha^{n-j},
 \qquad
 B_n(1-\alpha)=(-1)^nB_n(\alpha).
 \label{eq:Bernoulli-identities}
\end{equation}

\begin{theorem}[Shifted \(K\)-asymptotics]
\label{thm:shifted-forward}
Fix \(\alpha\in\C\).  As \(q\to\infty\) in a closed subsector of
\(|\arg q|<\pi\),
\begin{align}
 \log\Kfun(q+\alpha)
 \sim{}&\frac{q^2}{2}\log q-\frac{q^2}{4}
 +(\alpha-\tfrac12)q\log q
 +\frac{B_2(\alpha)}{2}\log q
 \notag\\
 &+\log A+\frac{\alpha^2-\alpha}{2}
 +\sum_{n\ge1}(-1)^{n+1}
   \frac{B_{n+2}(\alpha)}{n(n+1)(n+2)q^n}.
 \label{eq:shifted-forward-general}
\end{align}
After truncation at \(q^{-N}\), the remainder is \(O(q^{-N-1})\), with the
usual sectorial interpretation when a Stokes direction is approached.
\end{theorem}

\begin{proof}
Insert \(x=q+\alpha\) into \eqref{eq:unshifted-forward} and use
\[
 \log(q+\alpha)=\log q+
 \sum_{j\ge1}\frac{(-1)^{j+1}}{j}\frac{\alpha^j}{q^j},
 \qquad
 (q+\alpha)^{-2m}=q^{-2m}
 \sum_{j\ge0}(-1)^j\binom{2m+j-1}{j}\frac{\alpha^j}{q^j}.
\]
At each fixed inverse power only finitely many terms contribute.  The
polynomial and logarithmic part collects directly into the first line of
\eqref{eq:shifted-forward-general}.  For \(n\ge1\), the coefficient of
\(q^{-n}\) is a binomial convolution of Bernoulli numbers with powers of
\(\alpha\).  Identity \eqref{eq:Bernoulli-identities} reduces that convolution
to
\[
 (-1)^{n+1}\frac{B_{n+2}(\alpha)}{n(n+1)(n+2)}.
\]
This proves the formal expansion.  Because translation by a fixed \(\alpha\)
preserves closed proper subsectors and because the expansion
\eqref{eq:unshifted-forward} has a remainder after every finite truncation, the
coefficient collection also preserves the stated Poincar\'e remainder.
\end{proof}

\begin{remark}
Formula \eqref{eq:shifted-forward-general} is useful beyond the inverse
problem.  It packages every translated Euler--Maclaurin coefficient into a
single Bernoulli polynomial.  In particular, it makes the distinguished role
of the midpoint visible without inspecting coefficients one by one.
\end{remark}

\subsection{The unique logarithmic center and parity collapse}

The coefficient of \(q\log q\) in \eqref{eq:shifted-forward-general} vanishes
if and only if \(\alpha=1/2\).  At the same point,
\(B_{2m+1}(1/2)=0\) for all \(m\ge0\), by the reflection identity in
\eqref{eq:Bernoulli-identities}.  Thus the removal of the subleading logarithmic
monomial and the parity collapse of the tail occur at exactly the same center.

\begin{corollary}[Midpoint normal form]
\label{cor:midpoint-normal-form}
Let
\begin{equation}
 x=q+\frac12,
 \qquad
 \kappa=\log A-\frac18.
 \label{eq:q-kappa-def}
\end{equation}
Then
\begin{equation}
 \boxed{
 \log\Kfun\!\left(q+\frac12\right)
 \sim \Phi(q)-\frac1{24}\log q+\kappa
       +\sum_{m\ge1}c_mq^{-2m},}
 \label{eq:centered-forward}
\end{equation}
where
\begin{equation}
 \Phi(q)=\frac{q^2}{2}\log q-\frac{q^2}{4},
 \label{eq:Phi-def}
\end{equation}
\begin{equation}
 c_m=-\frac{B_{2m+2}(1/2)}{(2m)(2m+1)(2m+2)}
 =-\frac{(2^{-2m-1}-1)B_{2m+2}}
          {(2m)(2m+1)(2m+2)}.
 \label{eq:cm-def}
\end{equation}
Explicitly,
\begin{align}
 \log\Kfun\!\left(q+\frac12\right)
 \sim{}&\frac{q^2}{2}\log q-\frac{q^2}{4}
 -\frac1{24}\log q+\log A-\frac18
 \notag\\
 &-\frac{7}{5760q^2}
 +\frac{31}{161280q^4}
 -\frac{127}{1290240q^6}
 +\frac{511}{4866048q^8}
 -\cdots.
 \label{eq:centered-first}
\end{align}
\end{corollary}

\begin{proof}
Set \(\alpha=1/2\) in \eqref{eq:shifted-forward-general}.  Since
\(B_2(1/2)=-1/12\) and \((\alpha^2-\alpha)/2=-1/8\), the elementary part has
the claimed form.  Odd inverse powers vanish because the relevant Bernoulli
polynomials have odd degree and are evaluated at the reflection center.
For even degree, the multiplication identity gives
\[
 B_n(1/2)=(2^{1-n}-1)B_n,
\]
which yields \eqref{eq:cm-def}.
\end{proof}

\begin{methodbox}
A translation should be chosen before inversion.  For a general forward
expansion beginning with
\[
 a x^p(\log x-\beta)+b x^{p-1}\log x+\cdots,
\]
the shift \(x=q+\alpha\) changes the coefficient of
\(q^{p-1}\log q\) to \(ap\alpha+b\).  The normal-form choice is therefore
\[
 \alpha=-\frac{b}{ap}.
\]
For the \(K\)-function, \(a=1/2\), \(p=2\), and \(b=-1/2\), giving
\(\alpha=1/2\).  In problems derived from Euler--Maclaurin expansions, the same
center frequently converts ordinary Bernoulli numbers into Bernoulli
polynomials at a symmetry point and removes half of the power lattice.
\end{methodbox}

\section{The ordinary Lambert anchor}
\label{sec:ordinary-anchor}

\subsection{Exact inversion of the dominant core}

We first invert only \(\Phi(q)\) from \eqref{eq:Phi-def}.  Write
\begin{equation}
 \Phi(Q)=Y.
 \label{eq:Phi-anchor-equation}
\end{equation}
If \(w=2\log Q-1\), then \(Q^2=\e^{w+1}\) and
\[
 Y=\frac14w\e^{w+1}.
\]
Hence \(w\e^w=4Y/\e\), which proves the following elementary but central
identity.

\begin{proposition}[Lambert anchor]
\label{prop:Lambert-anchor}
For \(Y>0\), the positive solution of \(\Phi(Q)=Y\) is
\begin{equation}
 Q=\exp\!\left(\frac{1+\W_0(4Y/\e)}2\right)
 =2\sqrt{\frac{Y}{\W_0(4Y/\e)}}.
 \label{eq:Q-anchor}
\end{equation}
If
\begin{equation}
 w=\W_0(4Y/\e),\qquad L=\log Q=\frac{w+1}{2},
 \label{eq:w-L}
\end{equation}
then
\begin{equation}
 \Phi'(Q)=Q\log Q=QL.
 \label{eq:Phi-prime}
\end{equation}
\end{proposition}

The entire inverse problem is now a perturbation of an exactly inverted core.
Define
\begin{equation}
 R(q)=\kappa-\frac1{24}\log q+
      \sum_{m\ge1}c_mq^{-2m},
 \label{eq:R-ordinary}
\end{equation}
so that the centered equation is
\begin{equation}
 Y\sim\Phi(q)+R(q).
 \label{eq:centered-inverse-equation}
\end{equation}
At \(q=Q\), the leading perturbation is only \(O(\log Q)\), whereas
\(\Phi'(Q)=Q\log Q\).  The first displacement is consequently of order
\(Q^{-1}\), and the relative displacement is of order \(Q^{-2}\).

\subsection{A normalized implicit equation}

Put
\begin{equation}
 t=Q^{-2},\qquad q=QH,
 \qquad H=1+\sum_{n\ge1}a_n(L)t^n,
 \qquad w=2L-1,
 \label{eq:H-t-def}
\end{equation}
and retain \(\Delta=L/24-\kappa\).  Multiplying
\(\Phi(QH)+R(QH)-Y\) by \(4/Q^2\) gives the dimensionless equation
\begin{align}
 \calE(H;t,L):={}&H^2\bigl(w+2\log H\bigr)-w
 -4\Delta t-\frac{t}{6}\log H
 \notag\\
 &+4\sum_{m\ge1}c_m t^{m+1}H^{-2m}=0.
 \label{eq:E-ordinary}
\end{align}
This equation contains only nonnegative powers of \(t\) when \(H=1+O(t)\).
Moreover,
\begin{equation}
 \frac{\partial\calE}{\partial H}(1;0,L)=4L,
 \label{eq:E-linearization}
\end{equation}
which is invertible for large positive \(L\).  Coefficient matching is
therefore triangular.

\begin{theorem}[Triangular recurrence]
\label{thm:triangular-recurrence}
Let
\begin{equation}
 H_{n-1}(t)=1+\sum_{j=1}^{n-1}a_j(L)t^j.
 \label{eq:H-partial}
\end{equation}
Then the coefficients in \eqref{eq:main-inverse-series} are determined by
\begin{equation}
 \boxed{
 a_n(L)=-\frac1{4L}\coeff{t}{n}
 \calE\bigl(H_{n-1}(t);t,L\bigr),\qquad n\ge1.}
 \label{eq:an-recurrence}
\end{equation}
At order \(n\), only \(c_1,\ldots,c_{n-1}\) are needed.
Each \(a_n\) is a rational function of \(L\) with coefficients in
\(\Q[\kappa]\).  Moreover,
\[
 a_n(L)=O(1),\qquad L\to+\infty,
\]
so every \(a_n\) has a finite limit at infinity.
\end{theorem}

\begin{proof}
Assume \(a_1,\ldots,a_{n-1}\) have been chosen so that
\(\calE(H_{n-1};t,L)=O(t^n)\).  Replacing \(H_{n-1}\) by
\(H_{n-1}+a_nt^n\) changes the coefficient of \(t^n\) by
\[
 a_n\frac{\partial\calE}{\partial H}(1;0,L)=4La_n;
\]
all nonlinear effects of the new term begin at order \(t^{n+1}\).  Solving the
linear equation at order \(t^n\) gives \eqref{eq:an-recurrence}.  The assertion
about \(c_m\) follows from the factor \(t^{m+1}\) in
\eqref{eq:E-ordinary}.  Rationality follows inductively from closure of
\(\Q[\kappa,L,L^{-1}]\) under the finite coefficient operations involved.
Finally, \(\Delta=O(L)\).  If \(a_1,\ldots,a_{n-1}=O(1)\), then the
coefficient extracted in \eqref{eq:an-recurrence} is \(O(L)\): the only
large factor is the single affine occurrence of \(w=2L-1\), while all
Bernoulli-source terms are bounded.  Division by \(4L\) gives
\(a_n=O(1)\).  Since \(a_n\) is rational in \(L\), boundedness at
infinity implies the existence of a finite limit.
\end{proof}

\begin{algorithm}[Coefficient generation for the ordinary anchor]
\label{alg:ordinary-coefficients}
To compute \(a_1,\ldots,a_N\):
\begin{enumerate}
\item form \(c_1,\ldots,c_{N-1}\) from \eqref{eq:cm-def};
\item initialize \(H\leftarrow1\);
\item for \(n=1,\ldots,N\), expand \(\calE(H;t,L)\) only through \(t^n\),
      set \(a_n=-(4L)^{-1}[t^n]\calE\), and replace
      \(H\leftarrow H+a_nt^n\);
\item substitute \(t=Q^{-2}\) and add the outer shift \(1/2\).
\end{enumerate}
The truncation discipline is important: carrying terms above the current order
causes expression swell but contributes nothing to the next coefficient.
\end{algorithm}

\subsection{Explicit coefficients}

Besides \eqref{eq:a1}--\eqref{eq:a3}, the next coefficient is
\begin{align}
 a_4(L)={}&
 -\frac{(15L^3+34L^2+35L+15)\Delta^4}{24L^7}
 +\frac{(23L^2+32L+15)\Delta^3}{144L^6}
 \notag\\
 &+\frac{(105L^2-13L-39)\Delta^2}{11520L^5}
 +\frac{(930L^2-206L-77)\Delta}{967680L^4}
 \notag\\
 &+\frac{45720L^2-5435L+637}{464486400L^3}.
 \label{eq:a4}
\end{align}
The limits as \(L\to\infty\) begin
\begin{equation}
 a_1\to\frac1{24},\qquad
 a_2\to-\frac1{1152},\qquad
 a_3\to\frac1{27648},\qquad
 a_4\to-\frac5{2654208}.
 \label{eq:an-limits}
\end{equation}
These constants arise from repeated nonlinear interaction of the single
logarithmic defect \(-\tfrac1{24}\log q+\kappa\).  The Bernoulli tail begins to
contribute to \(a_2\).

\begin{corollary}[First practical approximants]
\label{cor:ordinary-practical}
With the notation of \eqref{eq:main-anchor-data}, define
\begin{align*}
 X_0(y)&=\frac12+Q,\\
 X_1(y)&=X_0(y)+\frac{a_1(L)}{Q},\\
 X_2(y)&=X_1(y)+\frac{a_2(L)}{Q^3},\\
 X_3(y)&=X_2(y)+\frac{a_3(L)}{Q^5}.
\end{align*}
Then
\begin{equation}
 \Kfun_+^{-1}(y)-X_N(y)=\BigO\!\left(Q^{-2N-1}\right),
 \qquad N=0,1,2,3.
 \label{eq:ordinary-practical-error}
\end{equation}
\end{corollary}

\section{Anchored Lagrange--B\"urmann inversion}
\label{sec:anchored-LB}

The recurrence above is efficient, but a coordinate-free formula explains why
it works and extends immediately to other power--logarithmic problems.  We use
the standard Lagrange--B\"urmann theorem~\cite{Henrici1974,FlajoletSedgewick2009}
in an anchor-adapted coordinate.

\begin{theorem}[Anchored Lagrange--B\"urmann formula]
\label{thm:anchored-LB}
Let
\begin{equation}
 F(x)=F_0(x)+R(x)
 \label{eq:F-split}
\end{equation}
be formal, analytic, or asymptotic data in a region where \(F_0\) is locally
invertible.  Put
\begin{equation}
 X=F_0^{-1}(Y),
 \qquad
 \calD_X=\frac1{F_0'(X)}\frac{\dd}{\dd X}.
 \label{eq:anchored-D}
\end{equation}
Whenever the resulting series is meaningful in the chosen formal topology or
asymptotic scale, the branch of \(F^{-1}\) asymptotic to \(X\) is
\begin{equation}
 \boxed{
 F^{-1}(Y)=X+
 \sum_{n\ge1}\frac{(-1)^n}{n!}\calD_X^{\,n-1}
 \left(\frac{R(X)^n}{F_0'(X)}\right).}
 \label{eq:anchored-LB}
\end{equation}
\end{theorem}

\begin{proof}
Let \(\chi=F_0^{-1}\), set \(u=F_0(x)\), and define
\[
 h(u)=R(\chi(u)).
\]
The equation \(Y=F(x)\) becomes \(Y=u+h(u)\).  The
Lagrange--B\"urmann formula for an arbitrary test function \(\varphi\) gives
\[
 \varphi(u)=\varphi(Y)+\sum_{n\ge1}\frac{(-1)^n}{n!}
 \left(\frac{\dd}{\dd Y}\right)^{n-1}
 \bigl(\varphi'(Y)h(Y)^n\bigr).
\]
Choose \(\varphi=\chi\).  Then
\[
 \varphi(Y)=X,
 \qquad
 \varphi'(Y)=\frac1{F_0'(X)},
 \qquad
 h(Y)=R(X),
 \qquad
 \frac{\dd}{\dd Y}=\frac1{F_0'(X)}\frac{\dd}{\dd X}=\calD_X.
\]
Substitution yields \eqref{eq:anchored-LB}.
\end{proof}

The first two correction blocks are worth recording explicitly:
\begin{align}
 F^{-1}(Y)={}&X-\frac{R}{F_0'}
 +\frac{RR'}{(F_0')^2}
 -\frac{F_0''R^2}{2(F_0')^3}+\cdots,
 \label{eq:anchored-first-two}
\end{align}
where every function on the right is evaluated at \(X\).  This formula is a
useful independent check on any coefficient recurrence.

For the ordinary \(K\)-anchor,
\begin{equation}
 F_0=\Phi,
 \qquad F_0'(Q)=Q\log Q,
 \qquad
 \calD_Q=\frac1{Q\log Q}\frac{\dd}{\dd Q},
 \label{eq:K-D-ordinary}
\end{equation}
and \eqref{eq:anchored-LB} becomes
\begin{equation}
 q(Y)\sim Q+
 \sum_{n\ge1}\frac{(-1)^n}{n!}
 \left(\frac1{Q\log Q}\frac{\dd}{\dd Q}\right)^{n-1}
 \left(\frac{R(Q)^n}{Q\log Q}\right),
 \label{eq:K-operator-series}
\end{equation}
with \(R\) interpreted through its asymptotic expansion
\eqref{eq:R-ordinary}.  Collecting equal powers of \(Q^{-1}\) in
\eqref{eq:K-operator-series} reproduces \eqref{eq:main-inverse-series} and the
recurrence \eqref{eq:an-recurrence}.

\begin{remark}[Why anchoring matters]
Applying the near-identity reversion formula directly to \(F\) is impossible
because \(F(x)-x\) is not small.  The exact inverse of the dominant core first
moves the problem into the coordinate \(u=F_0(x)\), where the perturbation is
small relative to the identity.  The operator \(\calD_X\) is simply
\(\dd/\dd u\) written back in the physical coordinate \(X\).
\end{remark}

\section{A refined \texorpdfstring{\(r\)-Lambert}{r-Lambert} anchor}
\label{sec:refined-anchor}

\subsection{An exactly invertible affine logarithmic core}

The ordinary anchor leaves the slowly varying term
\(-\tfrac1{24}\log q+\kappa\) to perturbation theory.  It can instead be
absorbed into the core.  The required inverse belongs to the \(r\)-Lambert
family introduced and studied by Mez\H{o} and Baricz~\cite{MezoBaricz2017}.

\begin{definition}[The \(r\)-Lambert function]
For fixed \(r\), let \(\rW_r(z)\) denote a branch of the inverse of
\begin{equation}
 s\longmapsto s\e^s+rs.
 \label{eq:rW-definition}
\end{equation}
Thus
\begin{equation}
 \rW_r(z)\e^{\rW_r(z)}+r\rW_r(z)=z.
 \label{eq:rW-equation}
\end{equation}
When \(r=0\), this reduces to the ordinary Lambert function.
\end{definition}

The following general observation is useful for many special-function
inversions.

\begin{theorem}[Exact inversion of an affine logarithmic power core]
\label{thm:affine-log-core}
Let
\begin{equation}
 \Psi(x)=a x^p\log x+b x^p+c\log x+d,
 \qquad ap\ne0.
 \label{eq:general-affine-core}
\end{equation}
Choose compatible branches and set
\begin{align}
 r&=\frac{c}{a}\exp\!\left(\frac{pb}{a}\right),
 \label{eq:general-r}\displaybreak[0]\\
 z&=\frac{p}{a}\exp\!\left(\frac{pb}{a}\right)
     \left(Y-d+\frac{cb}{a}\right).
 \label{eq:general-z}
\end{align}
If \(s=\rW_r(z)\), then a solution of \(\Psi(x)=Y\) is
\begin{equation}
 \boxed{x=\exp\!\left(\frac{s}{p}-\frac{b}{a}\right).}
 \label{eq:general-affine-inverse}
\end{equation}
Every local inverse branch is obtained by compatible choices of the logarithm
and of a branch of \(\rW_r\).
\end{theorem}

\begin{proof}
Put
\[
 s=p\log x+\frac{pb}{a}.
\]
Then \(\log x=s/p-b/a\) and \(x^p=\exp(s-pb/a)\).  The two power terms
collapse:
\[
 ax^p\log x+bx^p
 =\frac{a}{p}\exp\!\left(-\frac{pb}{a}\right)s\e^s.
\]
The remaining logarithmic term is \(cs/p-cb/a\).  Therefore \(\Psi(x)=Y\)
is equivalent to
\[
 s\e^s+\frac{c}{a}\exp\!\left(\frac{pb}{a}\right)s
 =\frac{p}{a}\exp\!\left(\frac{pb}{a}\right)
  \left(Y-d+\frac{cb}{a}\right),
\]
which is precisely \eqref{eq:rW-equation} with the data
\eqref{eq:general-r}--\eqref{eq:general-z}.
\end{proof}

\subsection{Application to the \texorpdfstring{\(K\)-function}{K-function}}

Define the refined core
\begin{equation}
 \Psi(q)=\frac{q^2}{2}\log q-\frac{q^2}{4}
 -\frac1{24}\log q+\kappa.
 \label{eq:Psi-K}
\end{equation}
In \cref{thm:affine-log-core}, take
\[
 a=\frac12,
 \qquad p=2,
 \qquad b=-\frac14,
 \qquad c=-\frac1{24},
 \qquad d=\kappa.
\]
Then
\begin{equation}
 r=-\frac1{12\e},
 \qquad
 z=\frac{48(Y-\kappa)+1}{12\e}.
 \label{eq:K-r-z}
\end{equation}
The positive large solution is exactly \(Q_*\) from
\eqref{eq:refined-anchor-data}.  Its derivative is
\begin{equation}
 \Psi'(q)=q\log q-\frac1{24q}.
 \label{eq:Psi-prime}
\end{equation}
The residual perturbation is now only
\begin{equation}
 S(q)=\sum_{m\ge1}c_mq^{-2m}.
 \label{eq:S-tail}
\end{equation}
Thus the first inverse displacement is
\[
 -\frac{S(Q_*)}{\Psi'(Q_*)}
 =\frac{7}{5760L_*Q_*^3}+O(Q_*^{-5}L_*^{-1}),
\]
which explains the improved scale before any detailed calculation.

\subsection{Refined recurrence and coefficients}

Let
\begin{equation}
 t=Q_*^{-2},\qquad q=Q_*H,
 \qquad
 H=1+\sum_{n\ge1}b_n(L_*)t^{n+1},
 \qquad w=2L_*-1.
 \label{eq:Hstar-def}
\end{equation}
The normalized equation is
\begin{align}
 \calE_*(H;t,L_*):={}&H^2\bigl(w+2\log H\bigr)-w
 -\frac{t}{6}\log H
 \notag\\
 &+4\sum_{m\ge1}c_m t^{m+1}H^{-2m}=0.
 \label{eq:E-refined}
\end{align}
The term \(-4\Delta t\) has disappeared because it belongs to the exact core.

\begin{theorem}[Refined triangular recurrence]
\label{thm:refined-recurrence}
Let
\begin{equation}
 H_{n-1}^*(t)=1+\sum_{j=1}^{n-1}b_j(L)t^{j+1}.
 \label{eq:Hstar-partial}
\end{equation}
Then
\begin{equation}
 \boxed{
 b_n(L)=-\frac1{4L}\coeff{t}{n+1}
 \calE_*\bigl(H_{n-1}^*(t);t,L\bigr),\qquad n\ge1.}
 \label{eq:bn-recurrence}
\end{equation}
Each \(b_n(L)\) is rational in \(L\), satisfies \(b_n(L)=O(L^{-1})\), and
uses only \(c_1,\ldots,c_n\).
\end{theorem}

\begin{proof}
The proof is identical to that of \cref{thm:triangular-recurrence}; the
linearization remains \(4L\).  The source term at the first unmatched order has
no factor growing faster than a fixed power of \(L^0\), while division by the
linearization contributes \(L^{-1}\).  Induction gives \(b_n=O(L^{-1})\).
\end{proof}

In addition to \eqref{eq:b1}--\eqref{eq:b3}, one finds
\begin{align}
 b_4(L)={}&-
 \frac{12877200L^3-703428L^2+57552L+539}
      {122624409600L^4},
 \label{eq:b4}\displaybreak[0]\\
 b_5(L)={}&
 \frac{1539856241280L^4-45140228640L^3+1666058823L^2
       -38145107L-25162137}
      {8034351316992000L^5}.
 \label{eq:b5}
\end{align}
Their leading large-\(L\) behavior recovers the negatives of the forward tail
coefficients divided by \(L\):
\begin{equation}
 Lb_1\to\frac7{5760},\quad
 Lb_2\to-\frac{31}{161280},\quad
 Lb_3\to\frac{127}{1290240},\quad
 Lb_4\to-\frac{511}{4866048}.
 \label{eq:bn-limits}
\end{equation}
This is exactly what the linear approximation
\(-S(Q_*)/\Psi'(Q_*)\) predicts.

The anchored operator form is
\begin{equation}
 q(Y)\sim Q_*+
 \sum_{n\ge1}\frac{(-1)^n}{n!}
 \left(\frac1{\Psi'(Q_*)}\frac{\dd}{\dd Q_*}\right)^{n-1}
 \left(\frac{S(Q_*)^n}{\Psi'(Q_*)}\right).
 \label{eq:refined-operator}
\end{equation}
For every fixed \(N\),
\begin{equation}
 \Kfun_+^{-1}(y)
 =\frac12+Q_*+\sum_{n=1}^{N}\frac{b_n(L_*)}{Q_*^{2n+1}}
 +\BigO\!\left(\frac{Q_*^{-2N-3}}{L_*}\right).
 \label{eq:refined-remainder}
\end{equation}

\begin{summarybox}
The ordinary anchor is simpler because it uses the ubiquitous \(\W_0\) and
already gives excellent approximations.  The refined anchor is structurally
cleaner: it absorbs every term of forward power gap two, leaves only the
Bernoulli tail, and gives a remainder smaller by two powers of the anchor.  The
price is evaluation of a single \(r\)-Lambert function, which itself is easily
computed by Newton iteration.
\end{summarybox}

\section{Expansion in elementary logarithms}
\label{sec:pure-log}

The Lambert anchors are the most compact and numerically stable form of the
answer.  For comparison with conventional asymptotic notation, one may expand
them further into nested logarithms.  This operation should be performed only
after the inverse has been organized into its distinct transseries sectors.

\subsection{Expansion of the ordinary anchor}

Set
\begin{equation}
 M=\log\!\left(\frac{4Y}{\e}\right),
 \qquad N=\log M.
 \label{eq:M-N-def}
\end{equation}
The classical principal-branch expansion~\cite{Corless1996} is
\begin{align}
 \W_0(4Y/\e)
 \sim{}&M-N+\frac{N}{M}
 +\frac{N(N-2)}{2M^2}
 +\frac{N(2N^2-9N+6)}{6M^3}
 \notag\\
 &+\frac{N(3N^3-22N^2+36N-12)}{12M^4}
 +\cdots.
 \label{eq:W-log-expansion}
\end{align}
Substitution into \(Q=2\sqrt{Y/\W_0(4Y/\e)}\) yields
\begin{align}
 Q\sim 2\sqrt{\frac{Y}{M}}\Bigg[{}&1+\frac{N}{2M}
 +\frac{N(3N-4)}{8M^2}
 +\frac{N(5N^2-16N+8)}{16M^3}
 \notag\\
 &+\frac{N(105N^3-568N^2+720N-192)}{384M^4}
 +\cdots\Bigg].
 \label{eq:Q-log-expansion}
\end{align}
The polynomials in \(N\) can be generated to arbitrary order either from the
standard Lambert polynomials or by substituting an ansatz in
\(w+\log w=M\).

The first inverse-power sector begins with
\begin{equation}
 \frac{a_1(L)}{Q}
 =\sqrt{\frac{M}{Y}}
 \left[\frac1{48}
 -\frac{\kappa+N/96}{M}
 +\BigO\!\left(\frac{N^2}{M^2}\right)\right].
 \label{eq:first-sector-log}
\end{equation}
Combining \eqref{eq:Q-log-expansion} and
\eqref{eq:first-sector-log} gives the first three sectors explicitly:
\begin{align}
 \Kfun_+^{-1}(y)
 ={}&2\sqrt{\frac{Y}{M}}
 \left[1+\frac{N}{2M}
 +\frac{N(3N-4)}{8M^2}
 +\frac{N(5N^2-16N+8)}{16M^3}
 +\BigO\!\left(\frac{N^4}{M^4}\right)\right]
 \notag\\
 &+\frac12
 +\sqrt{\frac{M}{Y}}
 \left[\frac1{48}
 -\frac{\kappa+N/96}{M}
 +\BigO\!\left(\frac{N^2}{M^2}\right)\right]
 +\BigO\!\left(\frac{M^{3/2}}{Y^{3/2}}\right).
 \label{eq:full-three-sector-log}
\end{align}

\subsection{Why the constant is beyond all logarithmic orders}

The leading sector has size \(Q\asymp\sqrt{Y/M}\).  Relative to it, the
centering term has size
\[
 \frac{1/2}{Q}\asymp\sqrt{\frac{M}{Y}}.
\]
Since \(Y=(\e/4)e^M\), this ratio is exponentially small as a function of the
coarse variable \(M\).  Hence no finite or infinite Poincar\'e series in powers
of \(1/M\) multiplying the leading scale can produce the constant \(1/2\).
The same is true, a fortiori, for the inverse-power sectors.

\begin{warningbox}
The statement
\[
 \Kfun_+^{-1}(y)\sim2\sqrt{Y/M}
 \left(1+\frac{N}{2M}+\cdots\right)
\]
is correct in the one-sector logarithmic sense, but it does not determine the
additive constant and therefore cannot be differentiated, composed, or used as
a high-accuracy numerical formula beyond its declared scale.  The Lambert-
anchored transseries \eqref{eq:main-inverse-series} is the more informative
object.
\end{warningbox}

\subsection{The transmonomial grid}

Expanding the ordinary anchor and every coefficient function produces a grid
whose typical monomials are
\begin{equation}
 Y^{1/2-n}M^{n-1/2-j}N^k,
 \qquad n,j\in\N_0,
 \label{eq:K-transmonomial-grid}
\end{equation}
with a degree bound on \(k\) depending linearly on \(j\), together with the
separate constant sector \(1/2\).  The ordering is lexicographic by exponential
size first:
\[
 \sqrt{\frac{Y}{M}}
 \succ 1
 \succ \sqrt{\frac{M}{Y}}
 \succ \left(\frac{M}{Y}\right)^{3/2}
 \succ\cdots,
\]
and inside each sector by descending powers of \(M\), with powers of
\(N=\log M\) treated as coefficients of slower growth.

This grid is not arbitrary.  In the centered forward expansion, the set of
power gaps is \(\{2,4,6,\ldots\}\).  The general support-transfer theorem in
\cref{thm:support-transfer} maps a total forward gap \(2n\) to inverse power
\(Q^{1-2n}\), exactly the odd inverse-power lattice in
\eqref{eq:main-inverse-series}.

\section{General inversion of power--logarithmic transseries}
\label{sec:general-theory}

We now isolate the structure that does not depend on the special coefficients
of the \(K\)-function.

\subsection{The universal Lambert anchor}

Consider the dominant model
\begin{equation}
 F_0(x)=a x^p(\log x-\beta),
 \qquad a\ne0,\quad p>0.
 \label{eq:general-F0}
\end{equation}
For the positive real problem assume \(a>0\) and take \(Y>0\) large.  Define
\begin{equation}
 U=\W_0\!\left(\frac{p}{a}e^{-p\beta}Y\right).
 \label{eq:general-U}
\end{equation}
Then
\begin{equation}
 X=\exp\!\left(\beta+\frac{U}{p}\right)
 =\left(\frac{pY}{aU}\right)^{1/p}
 \label{eq:general-X}
\end{equation}
solves \(F_0(X)=Y\).  Indeed,
\begin{equation}
 U=p(\log X-\beta),
 \qquad
 F_0'(X)=aX^{p-1}(U+1).
 \label{eq:general-anchor-identities}
\end{equation}

\begin{theorem}[Universal power--log anchor]
\label{thm:universal-anchor}
Every inverse problem with dominant term \eqref{eq:general-F0} admits an exact
Lambert anchor \eqref{eq:general-X}.  On a complex sector, replace \(\W_0\) by
a compatible branch \(\W_k\) and use the chosen logarithm in
\(X=\exp(\beta+U/p)\).  All algebraic corrections can then be generated by the
anchored formula \eqref{eq:anchored-LB}.
\end{theorem}

The examples in \cref{tab:anchor-examples} show that the same core occurs in
many classical inversions.

\begin{table}[htbp]
\centering
\caption{Typical power--logarithmic anchors.  Lower-order terms determine the
subsequent transseries but not the leading Lambert coordinate.}
\label{tab:anchor-examples}
\begin{tabularx}{\textwidth}{@{}lccc>{\raggedright\arraybackslash}X@{}}
\toprule
Problem & \(p\) & \(a\) & \(\beta\) & Dominant logarithm \\
\midrule
Inverse \(\Gamma\) & 1 & 1 & 1 &
 \(\log\Gamma(x)\sim x(\log x-1)\) \\
Inverse Barnes \(G\) & 2 & \(1/2\) & \(3/2\) &
 \(\log G(x+1)\sim\tfrac12x^2(\log x-\tfrac32)\) \\
Inverse centered \(K\) & 2 & \(1/2\) & \(1/2\) &
 \(\log K(q+1/2)\sim\tfrac12q^2(\log q-\tfrac12)\) \\
Generalized hyperfactorial with exponent \(k^r\) & \(r+1\) & \(1/(r+1)\) & \(1/(r+1)\) &
 \(\sum_{k\le n}k^r\log k\sim\frac{n^{r+1}}{r+1}
 (\log n-\frac1{r+1})\) \\
\bottomrule
\end{tabularx}
\end{table}

\subsection{The anchor differential operator}

Let \(U=p(\log X-\beta)\).  The operator in
\eqref{eq:anchored-D} takes the simple form
\begin{equation}
 \calD_X
 =\frac{X^{-p}}{a(U+1)}\left(X\frac{\dd}{\dd X}\right).
 \label{eq:general-D}
\end{equation}
For any exponent \(\mu\) and coefficient function \(f(U)\),
\begin{equation}
 \calD_X\bigl(X^\mu f(U)\bigr)
 =X^{\mu-p}\mathfrak{T}_\mu f(U),
 \label{eq:D-on-monomial}
\end{equation}
where
\begin{equation}
 \mathfrak{T}_\mu f
 :=\frac{\mu f+p f'}{a(U+1)}.
 \label{eq:Tmu-def}
\end{equation}
Thus every application of \(\calD_X\) lowers the power of \(X\) by exactly
\(p\), while acting on the slow coefficient through a first-order rational
differential operator.  This separation of fast and slow variables is the key
to an explicit support theorem.

\subsection{Grid-based perturbations}

Let \(\Lambda\subset(0,\infty)\) be a locally finite set of positive power
gaps.  Consider
\begin{equation}
 F(x)\sim F_0(x)+R(x),
 \qquad
 R(x)\sim\sum_{\lambda\in\Lambda}
 x^{p-\lambda}r_\lambda(U),
 \qquad U=p(\log x-\beta).
 \label{eq:general-R-grid}
\end{equation}
The coefficient functions may be rational functions of \(U\), polynomials in
\(\log U\) with rational-function coefficients, or elements of any
differential algebra \(\calA\) closed under \(f\mapsto f'\) and division by
\(U+1\).  Write \(\angles{\Lambda}\) for the additive semigroup generated by
\(\Lambda\).

For \(m\ge0\), define the iterated operator
\begin{equation}
 \mathfrak{T}^{\langle m\rangle}_{\mu}
 =\begin{cases}
 \operatorname{id},&m=0,\\[0.3em]
 \mathfrak{T}_{\mu-(m-1)p}\circ\cdots\circ
 \mathfrak{T}_{\mu-p}\circ\mathfrak{T}_{\mu},&m\ge1.
 \end{cases}
 \label{eq:T-iterated}
\end{equation}

\begin{theorem}[Support transfer and all-orders coefficients]
\label{thm:support-transfer}
Assume that every \(\sigma\in\angles{\Lambda}\) has only finitely many ordered
representations
\(\sigma=\lambda_1+\cdots+\lambda_n\) with
\(\lambda_j\in\Lambda\).  Then the inverse branch asymptotic to the Lambert
anchor \(X\) has the transseries
\begin{equation}
 \boxed{
 F^{-1}(Y)\sim X+
 \sum_{\sigma\in\angles{\Lambda}}X^{1-\sigma}C_\sigma(U),}
 \label{eq:general-inverse-grid}
\end{equation}
where
\begin{align}
 C_\sigma(U)
 ={}&\sum_{n\ge1}\frac{(-1)^n}{n!}
 \sum_{\substack{\lambda_1,\ldots,\lambda_n\in\Lambda\\
                   \lambda_1+\cdots+\lambda_n=\sigma}}
 \mathfrak{T}^{\langle n-1\rangle}_{(n-1)p+1-\sigma}
 \left(
   \frac{\prod_{j=1}^n r_{\lambda_j}(U)}{a(U+1)}
 \right).
 \label{eq:Csigma-formula}
\end{align}
For each fixed \(\sigma\), both sums are finite.  In particular, a forward
block with gap \(\lambda\) produces a first-order inverse block
\begin{equation}
 -X^{1-\lambda}\frac{r_\lambda(U)}{a(U+1)},
 \label{eq:first-gap-transfer}
\end{equation}
and an \(n\)-fold interaction with total gap \(\sigma\) always has fast power
\(X^{1-\sigma}\).
\end{theorem}

\begin{proof}
Expand the \(n\)-th summand of \eqref{eq:anchored-LB}.  An ordered product of
forward blocks indexed by \(\lambda_1,\ldots,\lambda_n\) contributes
\[
 R(X)^n\supset X^{np-\sigma}
 \prod_{j=1}^nr_{\lambda_j}(U),
 \qquad \sigma=\sum_{j=1}^n\lambda_j.
\]
After division by \(F_0'(X)=aX^{p-1}(U+1)\), its power is
\[
 \mu_0=(n-1)p+1-\sigma.
\]
Each of the remaining \(n-1\) applications of \(\calD_X\) lowers the power by
\(p\).  Formula \eqref{eq:D-on-monomial} therefore gives the final power
\[
 \mu_0-(n-1)p=1-\sigma
\]
and the slow coefficient
\[
 \mathfrak{T}^{\langle n-1\rangle}_{\mu_0}
 \left(\frac{\prod_jr_{\lambda_j}}{a(U+1)}\right).
\]
Multiplication by \((-1)^n/n!\) and summation over all ordered decompositions
of \(\sigma\) gives \eqref{eq:Csigma-formula}.  Local finiteness of the support
makes coefficient extraction legitimate.
\end{proof}

\begin{corollary}[Power-gap law]
\label{cor:power-gap-law}
If the smallest forward gap is \(\lambda_0>0\), then the relative correction to
the anchor begins at \(X^{-\lambda_0}\).  The \(n\)-fold self-interaction of
that smallest block begins at relative order \(X^{-n\lambda_0}\).  Mixed
interactions add their gaps.
\end{corollary}

\begin{proof}
Divide every term in \eqref{eq:general-inverse-grid} by \(X\).  The relative
power corresponding to total gap \(\sigma\) is \(X^{-\sigma}\).
\end{proof}

\subsection{The \texorpdfstring{\(K\)-function}{K-function} as a support theorem example}

For the ordinary centered \(K\)-anchor,
\begin{equation}
 p=2,
 \qquad a=\frac12,
 \qquad \beta=\frac12,
 \qquad U=2\log Q-1.
 \label{eq:K-general-parameters}
\end{equation}
The perturbation \eqref{eq:R-ordinary} can be written as
\begin{equation}
 R(Q)=Q^{2-2}r_2(U)+
 \sum_{m\ge1}Q^{2-(2m+2)}r_{2m+2}(U),
 \label{eq:K-R-grid}
\end{equation}
with
\begin{equation}
 r_2(U)=\kappa-\frac{U+1}{48},
 \qquad
 r_{2m+2}(U)=c_m.
 \label{eq:K-r-lambda}
\end{equation}
Thus \(\Lambda=\{2,4,6,\ldots\}\), and
\(\angles{\Lambda}=\{2,4,6,\ldots\}\).  The support-transfer theorem gives
powers \(Q^{1-2n}\), precisely the odd inverse powers in
\eqref{eq:main-inverse-series}.  It also gives, without recurrence,
\begin{equation}
 a_n(L)=C_{2n}(2L-1).
 \label{eq:an-Csigma}
\end{equation}

In the refined organization, the entire gap-two block is absorbed into
\(\Psi\).  The remaining gaps are \(\{4,6,8,\ldots\}\), so the first inverse
power is \(Q_*^{1-4}=Q_*^{-3}\).  This is the structural reason for the gain in
\eqref{eq:refined-inverse-series}.

\begin{methodbox}
The support-transfer theorem is a transseries analogue of a Newton polygon.
Before computing a single coefficient, it predicts:
\begin{enumerate}
\item which inverse powers can occur;
\item where each forward block first appears in the inverse;
\item which mixed products can contribute at a requested order;
\item which terms may be absorbed into a better exact anchor.
\end{enumerate}
This information is especially valuable in computer algebra, where a naive
univariate expansion often generates many coefficients that must later cancel.
\end{methodbox}

\subsection{A normalized equation for arbitrary power gaps}
\label{subsec:general-normalized-equation}

The closed coefficient formula \eqref{eq:Csigma-formula} is conceptually useful,
but direct coefficient matching is usually faster.  Write
\begin{equation}
 x=XH,
 \qquad
 H=1+\sum_{\sigma\in\angles{\Lambda}}h_\sigma(U)X^{-\sigma}.
 \label{eq:general-H}
\end{equation}
Since \(p(\log(XH)-\beta)=U+p\log H\), multiplication of
\(F(XH)-F_0(X)\) by \(p/(aX^p)\) gives
\begin{align}
 \calE_{p,a}(H):={}&H^p\bigl(U+p\log H\bigr)-U
 \notag\\
 &+\frac{p}{a}\sum_{\lambda\in\Lambda}X^{-\lambda}
 H^{p-\lambda}r_\lambda\bigl(U+p\log H\bigr)=0.
 \label{eq:general-normalized-E}
\end{align}
Its linearization is
\begin{equation}
 \left.\frac{\partial\calE_{p,a}}{\partial H}\right|_{H=1,X^{-1}=0}
 =p(U+1).
 \label{eq:general-E-linearization}
\end{equation}

\begin{theorem}[General triangular recurrence]
\label{thm:general-triangular}
Order \(\angles{\Lambda}\) increasingly.  Suppose all
\(h_\tau\) with \(\tau<\sigma\) have been determined, and put
\[
 H_{<\sigma}=1+\sum_{\tau<\sigma}h_\tau(U)X^{-\tau}.
\]
Then
\begin{equation}
 \boxed{
 h_\sigma(U)=-\frac1{p(U+1)}
 [X^{-\sigma}]\calE_{p,a}(H_{<\sigma}).}
 \label{eq:general-h-recurrence}
\end{equation}
Moreover \(h_\sigma=C_\sigma\), where \(C_\sigma\) is given by
\eqref{eq:Csigma-formula}.
\end{theorem}

\begin{proof}
The coefficient \(h_\sigma X^{-\sigma}\) enters the normalized equation
linearly at its first occurrence, with multiplier \(p(U+1)\) by
\eqref{eq:general-E-linearization}.  Every nonlinear expression containing the
new coefficient has total gap greater than \(\sigma\).  This proves the
recurrence.  Both the recurrence and the anchored Lagrange formula construct
the unique formal solution with \(H=1+o(1)\), so their coefficients agree.
\end{proof}

For a one-dimensional grid \(\Lambda\subset d\N\), set \(t=X^{-d}\); then
\eqref{eq:general-h-recurrence} reduces to ordinary power-series coefficient
matching.  The \(K\)-function has \(d=2\), which is why the recurrences
\eqref{eq:an-recurrence} and \eqref{eq:bn-recurrence} are especially simple.

\subsection{Coefficient algebras and nested logarithms}

The slow variable \(U\) is itself a Lambert function of \(Y\).  Several useful
coefficient algebras are stable under the operator \(\mathfrak{T}_\mu\):
\begin{align*}
 &\C(U),\\
 &\C(U)[\log U],\\
 &\C((U^{-1}))[\log U],\\
 &\C((U^{-1}))((\log U)^{-1}),
\end{align*}
with the last expression understood as an iterated formal field when needed.
The reason is elementary: differentiation preserves these classes and
\(1/(U+1)\) has a geometric expansion in \(U^{-1}\).  Thus a forward
polynomial--logarithmic coefficient system remains polynomial--logarithmic
after inversion.

If a coefficient contains deeper iterated logarithms, for example
\(\log\log x\), rewrite it as a function of
\[
 U,
 \qquad \log U,
 \qquad \log\log U,
 \quad\ldots.
\]
The chain rule introduces only rational combinations of these slow variables.
The support theorem is unchanged because it tracks the fast powers separately
from the coefficient algebra.

\begin{proposition}[Closure under inversion]
\label{prop:coefficient-closure}
Let \(\calA\) be a differential algebra of germs or formal series in \(U\)
that is closed under multiplication and under division by \(U+1\).  If every
\(r_\lambda\in\calA\), then every inverse coefficient
\(C_\sigma\in\calA\).
\end{proposition}

\begin{proof}
Formula \eqref{eq:Csigma-formula} uses only finite sums, products,
derivatives, and division by \(a(U+1)\).  These operations preserve
\(\calA\).
\end{proof}

\subsection{From grids to well-based transseries}

The finite-decomposition hypothesis in \cref{thm:support-transfer} is automatic
for a finitely generated positive grid.  It can be weakened in the usual Hahn-
and transseries setting~\cite{vanDerHoeven2006,Edgar2010}.  Let the power
gaps take values in an ordered abelian group and suppose their support is
well-ordered in the direction of increasing smallness.  Neumann's lemma then
ensures that finite sums of support elements remain well-ordered and that each
coefficient in products receives only finitely many contributions under the
standard Hahn-series hypotheses.  The proof of
\cref{thm:support-transfer} carries over verbatim.

What must not be dropped is some form of local finiteness.  A set of positive
gaps accumulating at zero can destroy the existence of a smallest correction
and make the formal implicit equation non-triangular.  In analytic applications
this signals that the chosen anchor has not separated the scales sufficiently;
a continuous Mellin spectrum or an additional transmonomial may be required.

\subsection{Mixed exponential sectors}

Suppose the forward transseries also contains exponentially small blocks,
\begin{equation}
 R_{\exp}(x)=
 \sum_{\omega\in\Omega}e^{-\omega x^\rho}
 x^{\nu_\omega}r_\omega(U).
 \label{eq:exp-forward}
\end{equation}
The anchored formula remains valid in a transseries field containing these
monomials.  At first order, each exponential block transfers to
\begin{equation}
 \delta x_\omega
 \sim-\frac{e^{-\omega X^\rho}X^{\nu_\omega}r_\omega(U)}
           {aX^{p-1}(U+1)}.
 \label{eq:exp-first-inverse}
\end{equation}
Products of forward exponentials add their actions
\(\omega X^\rho\), while algebraic gaps add exactly as in
\cref{thm:support-transfer}.  This observation is the algebraic core of the
Stokes-transfer discussion in \cref{sec:complex}.

\section{Normal forms and anchor design}
\label{sec:normal-forms}

The quality of an inverse transseries depends less on the number of terms than
on the choice of the core that is inverted exactly.  This section systematizes
that choice.

\subsection{Translation normal form}

Assume
\begin{equation}
 F(x)=a x^p(\log x-\beta)+b x^{p-1}\log x
 +c x^{p-1}+O(x^{p-2}\log^d x).
 \label{eq:translation-input}
\end{equation}
Set \(x=q+\alpha\).  The coefficient of \(q^{p-1}\log q\) becomes
\(ap\alpha+b\).  Hence
\begin{equation}
 \alpha_*=-\frac{b}{ap}
 \label{eq:optimal-translation}
\end{equation}
removes the largest untranslated defect.  The remaining coefficient of
\(q^{p-1}\) may then be absorbed into a redefinition of \(\beta\), if desired.
For Euler--Maclaurin expansions, \(\alpha_*\) often coincides with a symmetry
point of the relevant Bernoulli polynomials.

\subsection{Absorbing affine logarithmic defects}

After translation, one often obtains
\begin{equation}
 F(q)=a q^p\log q+bq^p+c\log q+d+S(q).
 \label{eq:affine-plus-tail}
\end{equation}
Theorem \ref{thm:affine-log-core} shows that the first four terms are not merely
``approximately invertible'': they have an exact inverse in terms of
\(\rW_r\).  If \(S(q)=O(q^{-\eta})\), this raises the first inverse correction
from the scale created by \(c\log q+d\) to
\(O(q^{1-p-\eta}/\log q)\).

For the \(K\)-function, the ordinary Lambert core leaves a gap-two defect,
whereas the affine core leaves a gap-four defect.  This is why the refined
anchor gains two powers of \(Q\).

\subsection{A hierarchy of possible anchors}

There is no absolute ``best'' anchor.  Three levels are useful:
\begin{enumerate}
\item \textbf{Elementary power--log anchor.}  Use
      \(a x^p(\log x-\beta)\) and ordinary \(\W\).  This maximizes portability
      and usually gives a closed all-orders coefficient grid.
\item \textbf{Affine logarithmic anchor.}  Add \(c\log x+d\) and use
      \(\rW_r\).  This often removes the leading inverse sector.
\item \textbf{Numerical implicit anchor.}  Absorb several algebraic tail terms
      into a monotone core and solve it numerically.  The residual transseries
      becomes shorter, but the anchor is no longer a named closed-form
      function.
\end{enumerate}
The first two retain symbolic transparency.  The third can be optimal for
high-precision numerical work.

\begin{methodbox}
Choose the anchor by a cost model, not by aesthetics.  A useful objective is to
minimize
\[
 \frac{|R(X)|}{|XF_0'(X)|},
\]
the predicted relative first correction, subject to the cost and branch
complexity of evaluating \(F_0^{-1}\).  The ordinary and refined \(K\)-anchors
are two Pareto-optimal points in this tradeoff.
\end{methodbox}

\section{Algorithms for symbolic and numerical inversion}
\label{sec:algorithms}

\subsection{Four complementary algorithms}

The general theory supports four computational strategies.

\paragraph{Triangular coefficient matching.}
Use \eqref{eq:general-normalized-E} and
\eqref{eq:general-h-recurrence}.  This is normally the fastest way to obtain a
long symbolic expansion because every order is solved by one linear division.

\paragraph{Anchored Lagrange--B\"urmann blocks.}
Use \eqref{eq:anchored-LB} or \eqref{eq:Csigma-formula}.  This gives a closed
formula, proves support constraints, and is excellent for checking a
recurrence.  Repeated differentiation can cause expression swell at high
orders.

\paragraph{Asymptotic fixed-point iteration.}
Rewrite
\[
 x=F_0^{-1}(Y-R(x))
\]
and iterate from \(x_0=X\).  Truncating after every operation makes this an
efficient formal algorithm.  Without truncation, intermediate expressions
usually grow unnecessarily.

\paragraph{Newton polishing.}
Given any asymptotic approximation \(x_0\), set
\begin{equation}
 x_{j+1}=x_j-\frac{\log\Kfun(x_j)-Y}
 {\log\Gamma(x_j)+x_j-\tfrac12-\tfrac12\log(2\pi)}.
 \label{eq:K-Newton}
\end{equation}
One or two iterations recover nearly all available working precision when the
Lambert anchor is already in the asymptotic regime.

\subsection{A sparse coefficient data structure}

For a general grid, represent a term by a pair
\((\sigma,f(U))\), standing for \(X^{-\sigma}f(U)\) in the relative factor
\(H\).  Store terms in a map keyed by \(\sigma\).  Multiplication adds keys;
differentiation applies \eqref{eq:D-on-monomial}; truncation discards keys above
the requested total gap.  This representation prevents accidental expansion
of \(U=\W(cY)\) into nested logarithms before the final presentation stage.

\begin{algorithm}[Sparse transseries reversion]
\label{alg:sparse-reversion}
Given \(p,a,\beta\), a gap set \(\Lambda\), coefficient functions
\(r_\lambda\), and a cutoff \(\Sigma\):
\begin{enumerate}
\item form the exact anchor \(U\) and \(X\) from
      \eqref{eq:general-U}--\eqref{eq:general-X};
\item initialize the sparse relative series \(H=1\);
\item enumerate \(\sigma\in\angles{\Lambda}\) in increasing order up to
      \(\Sigma\);
\item compute only the coefficient of \(X^{-\sigma}\) in
      \eqref{eq:general-normalized-E};
\item divide by \(-p(U+1)\) and insert the resulting \(h_\sigma\) into \(H\);
\item return \(x=XH\), optionally expanding \(U\) into elementary logarithms.
\end{enumerate}
\end{algorithm}

\subsection{Residual-driven adaptive truncation}

A practical implementation need not predict the optimal number of asymptotic
terms.  Compute successive truncations \(x_N\), evaluate the residual
\begin{equation}
 E_N=F(x_N)-Y,
 \label{eq:residual-def}
\end{equation}
and stop when \(|E_N/F'(x_N)|\) ceases to decrease or falls below the target
absolute error.  This criterion works for convergent local inverse series and
for divergent Poincar\'e series truncated near their least term.

\section{Analytic validity and error transfer}
\label{sec:analytic-validity}

Formal transseries become useful analytic statements only after residuals are
controlled.  For the positive real \(K\)-branch, this step is particularly
simple because the logarithmic derivative is explicit and positive in the
large-argument region.

\subsection{A residual certificate}

\begin{lemma}[Real residual bound]
\label{lem:residual-bound}
Let \(F\) be continuously differentiable and strictly increasing on an interval
\(I\).  Let \(x\in I\) solve \(F(x)=Y\), and let \(\widehat{x}\in I\).  If
\begin{equation}
 m=\inf_{u\in I}F'(u)>0,
 \label{eq:m-lower}
\end{equation}
then
\begin{equation}
 |x-\widehat{x}|
 \le\frac{|F(\widehat{x})-Y|}{m}.
 \label{eq:residual-bound}
\end{equation}
If \(F''\) is controlled and the residual is small, then more precisely
\begin{equation}
 x-\widehat{x}
 =-\frac{F(\widehat{x})-Y}{F'(\widehat{x})}
 +\BigO\!\left(
   \frac{F''(\widehat{x})[F(\widehat{x})-Y]^2}
        {F'(\widehat{x})^3}
 \right).
 \label{eq:residual-Newton-expansion}
\end{equation}
\end{lemma}

\begin{proof}
The mean value theorem gives
\[
 F(\widehat{x})-F(x)=F'(\xi)(\widehat{x}-x)
\]
for some \(\xi\) between \(x\) and \(\widehat{x}\), proving
\eqref{eq:residual-bound}.  Taylor expansion about \(\widehat{x}\), followed
by one formal solution step for \(x-\widehat{x}\), gives
\eqref{eq:residual-Newton-expansion}.
\end{proof}

For the \(K\)-problem we apply this lemma to
\(F(x)=\log\Kfun(x)\).  By \eqref{eq:dlogK},
\begin{equation}
 F'(x)=\log\Gamma(x)+x-\frac12-\frac12\log(2\pi)
 \sim x\log x.
 \label{eq:Fprime-asymptotic}
\end{equation}
Therefore a forward residual of size \(E\) produces an inverse error of size
approximately \(E/(x\log x)\).

\subsection{General error-transfer principle}

\begin{proposition}[Transfer of a forward remainder]
\label{prop:error-transfer}
Let \(F_0(x)=a x^p(\log x-\beta)\) with \(a>0\), and let
\(X=F_0^{-1}(Y)\).  Suppose a forward term or remainder has size
\begin{equation}
 E(X)=\BigO\!\left(X^{p-\rho}L^d\right),
 \qquad L=\log X,
 \qquad \rho>0.
 \label{eq:forward-error-shape}
\end{equation}
Then its first-order effect on the inverse has size
\begin{equation}
 \delta X
 =\BigO\!\left(X^{1-\rho}L^{d-1}\right).
 \label{eq:inverse-error-shape}
\end{equation}
The relative inverse error is
\(O(X^{-\rho}L^{d-1})\).
\end{proposition}

\begin{proof}
The derivative of the core is
\(F_0'(X)=aX^{p-1}(p(\log X-\beta)+1)\asymp X^{p-1}L\).
Use the first anchored correction \(-E(X)/F_0'(X)\).
\end{proof}

This proposition is the analytic counterpart of the formal power-gap law.  It
also shows why logarithmic defects are mild: one factor of \(\log X\) is lost
when a forward residual is divided by the derivative.

\subsection{Remainder of the ordinary \texorpdfstring{\(K\)-expansion}{K-expansion}}

Let
\begin{equation}
 H_N=1+\sum_{n=1}^{N}a_n(L)t^n,
 \qquad t=Q^{-2}.
 \label{eq:HN-ordinary}
\end{equation}
The recurrence makes the normalized residual
\begin{equation}
 \calE(H_N;t,L)=O(Lt^{N+1}).
 \label{eq:E-ordinary-order}
\end{equation}
Multiplying back by \(Q^2/4\) gives
\begin{equation}
 \log\Kfun\!\left(\frac12+QH_N\right)-Y
 =O(LQ^{-2N}).
 \label{eq:ordinary-forward-residual}
\end{equation}
Since the derivative is asymptotic to \(QL\),
\cref{lem:residual-bound} gives
\begin{equation}
 q-QH_N=O(Q^{-2N-1}),
 \label{eq:ordinary-inverse-error-proof}
\end{equation}
which proves \eqref{eq:main-remainder}.  The omitted part of the forward
Bernoulli series has exactly the same or smaller order and is therefore covered
by the estimate.

\subsection{Remainder of the refined expansion}

For
\begin{equation}
 H_N^*=1+\sum_{n=1}^{N}b_n(L_*)t^{n+1},
 \qquad t=Q_*^{-2},
 \label{eq:HN-refined}
\end{equation}
the normalized residual is
\begin{equation}
 \calE_*(H_N^*;t,L_*)=O(t^{N+2}).
 \label{eq:E-refined-order}
\end{equation}
After multiplication by \(Q_*^2/4\), the forward residual is
\(O(Q_*^{-2N-2})\).  Division by the derivative
\(\Psi'(Q_*)\sim Q_*L_*\) yields
\begin{equation}
 q-Q_*H_N^*=O(Q_*^{-2N-3}L_*^{-1}),
 \label{eq:refined-inverse-error-proof}
\end{equation}
which is \eqref{eq:refined-remainder}.

\begin{proposition}[Leading behavior of the refined coefficients]
\label{prop:bn-leading}
For every fixed \(n\ge1\),
\begin{equation}
 b_n(L)=-\frac{c_n}{L}+O(L^{-2}),
 \qquad L\to+\infty.
 \label{eq:bn-leading}
\end{equation}
\end{proposition}

\begin{proof}
At gap \(2n+2\), the direct contribution of the forward term
\(c_nq^{-2n}\) to the normalized equation is \(4c_nt^{n+1}\).  Division by
the linearization \(4L\) gives \(-c_n/L\).  Every nonlinear interaction
contains at least one previously computed \(b_j=O(L^{-1})\), and is therefore
\(O(L^{-2})\) after the final division by \(L\).
\end{proof}

\section{Numerical validation}
\label{sec:numerics}

\subsection{Reference evaluation}

For high-precision tests it is convenient to avoid direct evaluation of a
possibly enormous \(\Kfun(x)\).  Instead use
\begin{equation}
 \log\Kfun(x)=(x-1)\log\Gamma(x)-\log\Gfun(x),
 \label{eq:numeric-logK}
\end{equation}
and solve \(\log\Kfun(x)=Y\).  The ordinary and refined anchors provide
excellent initial guesses, so one Newton step is usually enough to create a
reference value at the displayed precision.

\subsection{Error table}

\Cref{tab:numerical-errors} compares signed errors.  In the ordinary columns,
\(X_0=Q+1/2\), \(X_1=X_0+a_1/Q\), and
\(X_2=X_1+a_2/Q^3\).  In the refined columns,
\(X_0^*=Q_*+1/2\), \(X_1^*=X_0^*+b_1/Q_*^3\), and
\(X_2^*=X_1^*+b_2/Q_*^5\).

\begin{table}[htbp]
\centering
\caption{Signed approximation errors \(X-\Kfun_+^{-1}(e^Y)\).}
\label{tab:numerical-errors}
\small
\setlength{\tabcolsep}{3pt}
\begin{tabular}{@{}r rrr rrr@{}}
\toprule
\(Y\) & \(X_0-x\) & \(X_1-x\) & \(X_2-x\) &
\(X_0^*-x\) & \(X_1^*-x\) & \(X_2^*-x\)\\
\midrule
10
& \(9.162\times10^{-3}\) & \(1.935\times10^{-5}\) & \(1.347\times10^{-7}\)
& \(-9.058\times10^{-6}\) & \(5.793\times10^{-8}\) & \(-1.632\times10^{-9}\)\\
100
& \(1.070\times10^{-3}\) & \(-2.046\times10^{-7}\) & \(5.165\times10^{-10}\)
& \(-4.585\times10^{-7}\) & \(5.911\times10^{-10}\) & \(-2.978\times10^{-12}\)\\
1000
& \(-1.503\times10^{-4}\) & \(-2.134\times10^{-8}\) & \(4.241\times10^{-12}\)
& \(-1.923\times10^{-8}\) & \(3.898\times10^{-12}\) & \(-2.917\times10^{-15}\)\\
\(10^6\)
& \(-3.809\times10^{-5}\) & \(-2.524\times10^{-13}\) & \(6.128\times10^{-19}\)
& \(-9.592\times10^{-13}\) & \(4.268\times10^{-19}\) & \(-6.569\times10^{-25}\)\\
\bottomrule
\end{tabular}
\end{table}

The data illustrate several structural predictions:
\begin{itemize}
\item one ordinary correction improves the anchor by roughly two powers of
      \(Q\);
\item the refined anchor without any tail correction is already comparable to
      an ordinary expansion with one correction;
\item each refined correction again gains approximately two powers of \(Q_*\),
      with an additional slow factor consistent with \(L_*^{-1}\);
\item the asymptotic formulas are useful well before \(Y\) becomes enormous.
\end{itemize}

For reference, the corresponding exact inverse values are
\begin{align*}
 \Kfun_+^{-1}(e^{10})&=4.9661579454967092289\ldots,\\
 \Kfun_+^{-1}(e^{100})&=10.915136726993084682\ldots,\\
 \Kfun_+^{-1}(e^{1000})&=27.284589040600087008\ldots,\\
 \Kfun_+^{-1}(e^{10^6})&=584.23586441169240127\ldots.
\end{align*}

\begin{conjecture}[Enveloping refined expansion]
\label{conj:enveloping}
For sufficiently large positive \(Y\), the refined truncations in
\eqref{eq:refined-inverse-series} alternate around the exact inverse and the
error after truncation has the sign and approximate magnitude of the first
omitted term.  Equivalently, the refined series is eventually enveloping on the
positive real ray.
\end{conjecture}

The sign pattern in \cref{tab:numerical-errors} and the alternating Bernoulli
input support this conjecture, but a proof would require a sign-controlled
remainder representation for the centered Hurwitz-zeta derivative and a
monotone transfer through the inverse map.

\section{Complex branches and exponentially small sectors}
\label{sec:complex}

\subsection{Sectorial Lambert branches}

The forward expansion \eqref{eq:centered-forward} is valid in any closed
subsector of the cut plane once a branch of \(\log q\) is fixed.  The dominant
equation \(\Phi(Q)=Y\) has sectorial solutions
\begin{equation}
 w_k=\W_k\!\left(\frac{4Y}{\e}\right),
 \qquad
 Q_k=\exp\!\left(\frac{w_k+1}{2}\right),
 \qquad k\in\mathbb{Z},
 \label{eq:complex-Qk}
\end{equation}
provided the branch of \(\log Q_k\) is chosen so that
\(2\log Q_k-1=w_k\).  The algebraic inverse transseries on that sector is
obtained from \eqref{eq:main-inverse-series} by replacing \((Q,L)\) with
\((Q_k,(w_k+1)/2)\).  The coefficient functions themselves do not change.

For the general anchor \eqref{eq:general-F0}, the branch formula is
\begin{equation}
 U_k=\W_k\!\left(\frac{p}{a}e^{-p\beta}Y\right),
 \qquad
 X_k=\exp\!\left(\beta+\frac{U_k}{p}\right).
 \label{eq:general-complex-anchor}
\end{equation}
This expression is safer than writing a free \(p\)-th root: it records the
compatibility between the logarithm and the Lambert branch.  For integer
\(p\), different \(k\) asymptotically generate the expected root rotations,
but the correspondence is modified by the slowly varying logarithmic phase.

\subsection{Critical points and branch loss}

The anchor derivative vanishes when
\begin{equation}
 U+1=0.
 \label{eq:anchor-critical}
\end{equation}
In the ordinary Lambert coordinate this is the familiar branch point of
\(\W\) at argument \(-1/\e\).  Near such a point, division by \(U+1\) in
\eqref{eq:Tmu-def} is singular and the inverse is no longer described by an
ordinary power transseries in the same coordinate.  A square-root local model
or a uniform coalescing-saddle expansion is required.  The large positive
\(K\)-branch stays far from this obstruction because \(U\to+\infty\).

The refined \(r\)-Lambert anchor has analogous critical points determined by
\begin{equation}
 \e^s(1+s)+r=0.
 \label{eq:rW-critical}
\end{equation}
For \(r=-1/(12\e)\) and the large positive branch, the derivative is positive
and no ambiguity occurs.

\subsection{Stokes transfer to the inverse}

The algebraic expansion of \(\zeta'(-1,q+1/2)\), and equivalently of the Barnes
\(G\)-function, possesses exponentially improved refinements~\cite{Nemes2014,Nemes2017}.  Across Stokes
rays, exponentially small forward contributions switch on smoothly.  Suppose
schematically that
\begin{equation}
 F(q)=F_{\mathrm{alg}}(q)+\mathfrak{S}(q)e^{-\chi(q)},
 \qquad \Re\chi(q)>0,
 \label{eq:forward-Stokes}
\end{equation}
where \(F_{\mathrm{alg}}\) is the optimally truncated algebraic expansion.
Let \(q_{\mathrm{alg}}\) solve
\(F_{\mathrm{alg}}(q_{\mathrm{alg}})=Y\).  The first inverse Stokes sector is
then
\begin{equation}
 q-q_{\mathrm{alg}}
 \sim-\frac{\mathfrak{S}(q_{\mathrm{alg}})
 e^{-\chi(q_{\mathrm{alg}})}}
 {F_{\mathrm{alg}}'(q_{\mathrm{alg}})}.
 \label{eq:inverse-Stokes-first}
\end{equation}
Higher powers and interactions follow from the same anchored
Lagrange--B\"urmann formula.  Thus inversion does not create a new exponential
action at first order; it transports the forward action and divides its
amplitude by the logarithmic derivative.

\begin{statusbox}
Equation \eqref{eq:inverse-Stokes-first} is a rigorous formal transfer rule once
an exponentially improved forward expansion with controlled remainder is
supplied.  Determining the complete Stokes multipliers of the inverse
\(K\)-function on all complex sectors is a separate problem and is not claimed
here.  Existing remainder representations for the Hurwitz zeta derivative and
Barnes \(G\) provide the natural starting point.
\end{statusbox}

\section{Extensions to related special functions}
\label{sec:extensions}

\subsection{Generalized hyperfactorial products}

For a fixed real \(r>-1\), consider the discrete family
\begin{equation}
 H_r(n)=\prod_{k=1}^{n}k^{k^r},
 \qquad
 \log H_r(n)=\sum_{k=1}^{n}k^r\log k.
 \label{eq:generalized-hyperfactorial-product}
\end{equation}
Euler--Maclaurin summation gives the dominant term
\begin{equation}
 \log H_r(n)
 \sim\frac{n^{r+1}}{r+1}
 \left(\log n-\frac1{r+1}\right).
 \label{eq:Hr-dominant}
\end{equation}
Put \(p=r+1\).  The universal anchor is therefore
\begin{equation}
 X_r(Y)=
 \left(
   \frac{p^2Y}{\W_0(p^2Y/\e)}
 \right)^{1/p}.
 \label{eq:Hr-anchor}
\end{equation}
The ordinary hyperfactorial corresponds to \(r=1\), \(p=2\), and reproduces
\(Q=2\sqrt{Y/\W_0(4Y/\e)}\).  Once a chosen analytic continuation is expanded
in a translated Euler--Maclaurin normal form, the support-transfer theorem
generates its inverse transseries without further conceptual changes.

For integer \(r\), the coefficients are naturally expressed through Bernoulli
polynomials and generalized Glaisher--Kinkelin constants.  The translation
normal form predicts which center should be used before inversion; the center
must be derived from the coefficient of the \(n^r\log n\) boundary term rather
than assumed in advance.

\subsection{Inverse gamma and Barnes functions}

The inverse gamma and inverse Barnes-\(G\) problems are neighboring members of
the same class:
\begin{align*}
 \log\Gamma(x)&=x(\log x-1)-\frac12\log x
 +\frac12\log(2\pi)+O(x^{-1}),\\
 \log\Gfun(x+1)&=\frac{x^2}{2}
 \left(\log x-\frac32\right)
 -\frac1{12}\log x+\zeta'(-1)+O(x^{-2}).
\end{align*}
The ordinary anchors are respectively
\begin{align}
 X_\Gamma(Y)&=\frac{Y}{\W_0(Y/\e)},
 \label{eq:gamma-anchor}\displaybreak[0]\\
 X_G(Y)&=2\sqrt{\frac{Y}{\W_0(4Ye^{-3})}}.
 \label{eq:G-anchor}
\end{align}
The same affine-core theorem can absorb their logarithmic and constant defects
into appropriate \(r\)-Lambert anchors.  The difference among these functions
lies primarily in the translated support and coefficient data, not in the
inversion mechanism.

\subsection{Barnes multiple gamma functions}

After a conventional sign normalization, logarithms of Barnes multiple gamma
functions~\cite{Vardi1988} have dominant terms of the form \(x^p\log x\), with \(p\) equal to
the order of the relevant multiple-gamma level.  Their Euler--Maclaurin tails
are governed by multiple Bernoulli polynomials.  The general framework of
\cref{sec:general-theory} therefore predicts:
\begin{enumerate}
\item a Lambert anchor with the corresponding exponent \(p\);
\item a translation normal form determined by the first boundary logarithm;
\item an inverse support equal to the additive semigroup generated by the
      translated multiple-Bernoulli gaps;
\item coefficient functions in a differential algebra generated by the slow
      Lambert variable and the parameters of the multiple gamma function.
\end{enumerate}
Deriving explicit centered coefficients is largely an algebraic task once the
forward multiple-Bernoulli expansion is fixed.

\subsection{Parameter-dependent and double-scaling inversions}

If \(a,p,\beta\), or the forward coefficients depend on a parameter
\(\theta\), the preceding results remain uniform provided
\begin{equation}
 |a(\theta)|\ge a_0>0,
 \qquad
 |U(\theta)+1|\ge u_0>0,
 \label{eq:uniform-conditions}
\end{equation}
and the coefficient algebra has uniform asymptotic bounds.  New behavior occurs
when \(U+1\) approaches zero or when a forward gap tends to zero with
\(\theta\).  Such limits require a double-scaling variable and generally
replace the Lambert normal form by a uniform canonical function.

\section{Open problems and research directions}
\label{sec:research-directions}

\subsection{Exponentially improved inverse expansions}

Combine a terminant-based or resurgence-based expansion for
\(\zeta'(-1,z)\) with the anchored formula to compute explicit inverse Stokes
multipliers.  The leading transfer is \eqref{eq:inverse-Stokes-first}; the
nontrivial work is to organize interactions between the algebraic inverse
sectors and the exponentially small forward sectors near Stokes smoothing
regions.

\subsection{A proof of the enveloping property}

\Cref{conj:enveloping} should be approachable through a sign-definite integral
remainder for the centered expansion.  The midpoint shift may be essential:
its Bernoulli-polynomial coefficients alternate cleanly, while the unshifted
series mixes parity and obscures monotonicity.  A successful proof would give
rigorous one-sided bounds for \(\Kfun_+^{-1}\) from successive refined
truncations.

\subsection{Global real and complex branch census}

The two positive stationary points described in
\cref{sec:K-background} account for the local real branch changes, but the
analytic continuation has a much richer complex critical-value set.  A useful
program is to locate zeros of
\begin{equation}
 \log\Gamma(z)+z-\frac12-\frac12\log(2\pi)=0
 \label{eq:complex-critical-equation}
\end{equation}
and map them through \(\Kfun\).  These critical values determine the branch
points and the natural domains of inverse branches.  The Lambert branch model
\eqref{eq:complex-Qk} gives their large-index asymptotic organization.

\subsection{Automatic normal-form selection}

A symbolic system could inspect a forward transseries, choose a translation
that removes the largest subleading logarithmic monomial, test whether the
remaining affine logarithmic core is \(r\)-Lambert invertible, and then build
the sparse inverse support automatically.  The selection criterion in
\cref{sec:normal-forms} suggests a deterministic optimization problem over a
small library of exact anchors.

\subsection{Formal verification}

The algebraic core is well suited to mechanization.  A formal development can
separate:
\begin{enumerate}
\item the exact identities for \(\Kfun\), Hurwitz zeta, and Barnes \(G\);
\item finite coefficient identities in Bernoulli polynomials;
\item the formal Lagrange--B\"urmann theorem in a complete filtered ring;
\item the support-transfer theorem for finitely generated grids;
\item analytic remainder estimates on the positive ray.
\end{enumerate}
The triangular recurrences are especially friendly to proof assistants because
each coefficient theorem reduces to a finite ring identity after the previous
orders are available.

\subsection{Arithmetic of inverse coefficients}

The coefficients \(c_m\) have controlled Bernoulli denominators, while
\(a_n(L)\) and \(b_n(L)\) are rational functions with highly structured integer
numerators.  It is natural to ask for denominator bounds, \(p\)-adic valuations,
and combinatorial formulas analogous to the Stirling-number formulas for the
ordinary Lambert expansion.  The operator formula
\eqref{eq:Csigma-formula} provides a starting point: its repeated first-order
operators suggest Bell-polynomial and rooted-tree interpretations.

\section{Conclusion}
\label{sec:conclusion}

The large inverse of the generalized hyperfactorial is controlled by a
quadratic logarithmic saddle, but its clean transseries is visible only after
two normalizations.  First, passing from \(x\) to \(q=x-1/2\) removes the
subleading \(q\log q\) term and converts the Hurwitz-zeta tail into an even
Bernoulli-polynomial series.  Second, solving the dominant core exactly with a
Lambert function separates nested logarithms from genuinely new inverse-power
sectors.

The ordinary anchor
\[
 Q=2\sqrt{\frac{\log y}{\W_0(4\log y/\e)}}
\]
leads to
\[
 \Kfun_+^{-1}(y)
 \sim\frac12+Q+\frac{a_1}{Q}+\frac{a_2}{Q^3}+\cdots.
\]
The affine logarithmic core has an exact \(r\)-Lambert inverse and produces the
sharper organization
\[
 \Kfun_+^{-1}(y)
 \sim\frac12+Q_*+\frac{b_1}{Q_*^3}
 +\frac{b_2}{Q_*^5}+\cdots.
\]
Both expansions have explicit all-orders operator formulas, sparse triangular
recurrences, and direct residual estimates.

More generally, an anchor of shape \(a x^p(\log x-\beta)\) converts the inverse
problem into a filtered near-identity problem.  Forward power gaps add under
nonlinear interaction and transfer to inverse powers by the rule
\(\sigma\mapsto1-\sigma\).  This support law is the central reusable result: it
predicts the architecture of an inverse transseries before any coefficients are
computed, applies uniformly to gamma, Barnes, and generalized hyperfactorial
families, and provides a natural interface between symbolic algebra, numerical
Newton correction, and future formal verification.

\appendix

\section{Low-order coefficient audit}
\label{app:coefficient-audit}

This appendix records enough intermediate algebra to make the displayed
coefficients independently checkable.

\subsection{Forward coefficients}

From \eqref{eq:cm-def},
\begin{center}
\begin{tabular}{@{}ccl@{}}
\toprule
\(m\) & \(c_m\) & contribution to \(\log\Kfun(q+1/2)\)\\
\midrule
1 & \(-7/5760\) & \(-7/(5760q^2)\)\\
2 & \(31/161280\) & \(31/(161280q^4)\)\\
3 & \(-127/1290240\) & \(-127/(1290240q^6)\)\\
4 & \(511/4866048\) & \(511/(4866048q^8)\)\\
5 & \(-1414477/7380172800\) & \(-1414477/(7380172800q^{10})\)\\
\bottomrule
\end{tabular}
\end{center}

\subsection{Ordinary recurrence at the first orders}

Write
\[
 H=1+a_1t+a_2t^2+a_3t^3+O(t^4).
\]
Expanding \eqref{eq:E-ordinary} gives
\begin{align}
 [t]\calE={}&4(La_1-\Delta),
 \label{eq:audit-E1}\displaybreak[0]\\
 [t^2]\calE={}&
 2La_1^2+4La_2+2a_1^2-\frac{a_1}{6}-\frac7{1440},
 \label{eq:audit-E2}\displaybreak[0]\\
 [t^3]\calE={}&
 4La_1a_2+4La_3+\frac23a_1^3+\frac1{12}a_1^2
 +4a_1a_2+\frac7{720}a_1-\frac16a_2+\frac{31}{40320}.
 \label{eq:audit-E3}
\end{align}
The first equation gives \(a_1=\Delta/L\).  Substitution into the second gives
\eqref{eq:a2}; substitution of both into the third gives \eqref{eq:a3}.
These equations are also useful unit tests for a symbolic implementation.

\subsection{Refined recurrence at the first orders}

Write
\[
 H=1+b_1t^2+b_2t^3+b_3t^4+O(t^5).
\]
Expansion of \eqref{eq:E-refined} gives
\begin{align}
 [t^2]\calE_*={}&4Lb_1-\frac7{1440},
 \label{eq:audit-b1}\displaybreak[0]\\
 [t^3]\calE_*={}&4Lb_2-\frac16b_1+\frac{31}{40320},
 \label{eq:audit-b2}\displaybreak[0]\\
 [t^4]\calE_*={}&
 2Lb_1^2+4Lb_3+2b_1^2+\frac7{720}b_1
 -\frac16b_2-\frac{127}{322560}.
 \label{eq:audit-b3}
\end{align}
Solving successively reproduces \eqref{eq:b1}--\eqref{eq:b3}.

\subsection{Independent first-order check}

The first anchored correction is \(-R(Q)/\Phi'(Q)\).  Since
\[
 R(Q)=\kappa-\frac{L}{24}+\frac{c_1}{Q^2}+O(Q^{-4})
 =-\Delta+\frac{c_1}{Q^2}+O(Q^{-4}),
\]
we obtain
\[
 -\frac{R(Q)}{QL}
 =\frac{\Delta}{LQ}-\frac{c_1}{LQ^3}+O(Q^{-5}),
\]
which verifies \(a_1\) and the direct Bernoulli contribution
\(7/(5760L)\) inside \(a_2\).  The remaining pieces of \(a_2\) are the
quadratic terms in \eqref{eq:anchored-first-two}.

\section{Reference implementation}
\label{app:code}

The following self-contained Python code uses \texttt{mpmath}.  It accepts
\(Y=\log y\) rather than \(y\), avoiding overflow when the original argument
is enormous.  The functions \texttt{ordinary} and \texttt{refined} implement
the first three correction blocks displayed in the article; \texttt{polish}
uses the exact logarithmic derivative for optional Newton refinement.

\begin{lstlisting}
import mpmath as mp

mp.mp.dps = 80
KAPPA = mp.log(mp.glaisher) - mp.mpf(1) / 8


def log_k(x):
    """log K(x), evaluated through Gamma and Barnes G."""
    x = mp.mpf(x)
    return (x - 1) * mp.loggamma(x) - mp.log(mp.barnesg(x))


def d_log_k(x):
    """Exact derivative of log K(x)."""
    x = mp.mpf(x)
    return (mp.loggamma(x) + x - mp.mpf('0.5')
            - mp.log(2 * mp.pi) / 2)


def ordinary_coefficients(L):
    """a_1, a_2, a_3 for the ordinary Lambert anchor."""
    D = L / 24 - KAPPA
    a1 = D / L
    a2 = (mp.mpf(7) / (5760 * L)
          + D / (24 * L**2)
          - (L + 1) * D**2 / (2 * L**3))
    a3 = ((3 * L**2 + 5 * L + 3) * D**3 / (6 * L**5)
          - (4 * L + 3) * D**2 / (48 * L**4)
          + (1 - 7 * L) * D / (1920 * L**3)
          + (49 - 186 * L) / (967680 * L**2))
    return [a1, a2, a3]


def ordinary(Y, terms=3):
    """Asymptotic inverse of log K(x)=Y, ordinary anchor."""
    Y = mp.mpf(Y)
    if not 0 <= terms <= 3:
        raise ValueError("terms must be between 0 and 3")
    w = mp.lambertw(4 * Y / mp.e)
    Q = 2 * mp.sqrt(Y / w)
    L = mp.log(Q)
    x = Q + mp.mpf('0.5')
    for n, a in enumerate(ordinary_coefficients(L)[:terms], start=1):
        x += a / Q**(2 * n - 1)
    return mp.re(x)


def r_lambert_positive(z, r, max_steps=30):
    """Positive large solution of s*exp(s)+r*s=z by Newton iteration."""
    z = mp.mpf(z)
    r = mp.mpf(r)
    s = mp.lambertw(z)  # excellent initial value when z is large
    for _ in range(max_steps):
        es = mp.exp(s)
        f = s * es + r * s - z
        fp = es * (s + 1) + r
        step = f / fp
        s -= step
        if abs(step) <= mp.eps * max(1, abs(s)):
            break
    return s


def refined_coefficients(L):
    """b_1, b_2, b_3 for the refined r-Lambert anchor."""
    b1 = mp.mpf(7) / (5760 * L)
    b2 = -(186 * L - 49) / (967680 * L**2)
    b3 = ((45720 * L**2 - 5435 * L + 637)
          / (464486400 * L**3))
    return [b1, b2, b3]


def refined(Y, terms=3):
    """Asymptotic inverse of log K(x)=Y, refined r-Lambert anchor."""
    Y = mp.mpf(Y)
    if not 0 <= terms <= 3:
        raise ValueError("terms must be between 0 and 3")
    r = -1 / (12 * mp.e)
    z = (48 * (Y - KAPPA) + 1) / (12 * mp.e)
    s = r_lambert_positive(z, r)
    Q = mp.exp((s + 1) / 2)
    L = mp.log(Q)
    x = Q + mp.mpf('0.5')
    for n, b in enumerate(refined_coefficients(L)[:terms], start=1):
        x += b / Q**(2 * n + 1)
    return x


def polish(Y, x, steps=1):
    """Newton polishing using exact log K and its exact derivative."""
    Y = mp.mpf(Y)
    x = mp.mpf(x)
    for _ in range(steps):
        x -= (log_k(x) - Y) / d_log_k(x)
    return x


if __name__ == "__main__":
    for Y in [10, 100, 1000, mp.mpf(10)**6]:
        x0 = refined(Y, terms=2)
        x = polish(Y, x0, steps=2)
        print("Y =", Y)
        print("asymptotic =", mp.nstr(x0, 30))
        print("polished   =", mp.nstr(x, 30))
        print("residual   =", mp.nstr(log_k(x0) - Y, 8))
\end{lstlisting}

\section{Notation and branch conventions}
\label{app:notation}

\begin{longtable}{@{}p{0.20\textwidth}p{0.73\textwidth}@{}}
\toprule
Symbol & Meaning\\
\midrule
\endhead
\(\Kfun(z)\) & Generalized hyperfactorial normalized by
\(\Kfun(z+1)=z^z\Kfun(z)\), \(\Kfun(1)=1\).\\
\(\Kfun_+^{-1}\) & Large positive real inverse branch.\\
\(\Gfun(z)\) & Barnes \(G\)-function.\\
\(A\) & Glaisher--Kinkelin constant.\\
\(Y\) & \(\log y\), the logarithm of the argument of the inverse.\\
\(q\) & Centered variable \(x-1/2\).\\
\(w\) & \(\W_0(4Y/\e)\).\\
\(Q\) & Ordinary Lambert anchor \(\exp((w+1)/2)\).\\
\(L\) & \(\log Q=(w+1)/2\).\\
\(Q_*,L_*\) & Refined \(r\)-Lambert anchor and its logarithm.\\
\(\rW_r\) & Inverse of \(s\mapsto s\e^s+rs\).\\
\(B_n(x)\) & Bernoulli polynomial; \(B_n=B_n(0)\).\\
\(\Lambda\) & Set of positive forward power gaps.\\
\(\angles{\Lambda}\) & Additive semigroup generated by \(\Lambda\).\\
\(\calD_X\) & Anchor derivative \((F_0'(X))^{-1}\dd/\dd X\).\\
\bottomrule
\end{longtable}

\section*{Acknowledgment of computational checks}

The rational coefficient identities and the numerical table were independently
checked by symbolic series expansion and high-precision numerical evaluation of
\eqref{eq:numeric-logK}.  These checks are not substitutes for the proofs; they
serve as regression tests for signs, shifts, and denominator arithmetic.

\begin{thebibliography}{99}

\bibitem{Adamchik1998}
V.~S. Adamchik,
``Polygamma functions of negative order,''
\emph{Journal of Computational and Applied Mathematics} \textbf{100} (1998),
191--199.
\url{https://doi.org/10.1016/S0377-0427(98)00192-7}.

\bibitem{Corless1996}
R.~M. Corless, G.~H. Gonnet, D.~E.~G. Hare, D.~J. Jeffrey, and D.~E. Knuth,
``On the Lambert \(W\) function,''
\emph{Advances in Computational Mathematics} \textbf{5} (1996), 329--359.
\url{https://doi.org/10.1007/BF02124750}.

\bibitem{DLMF5}
NIST Digital Library of Mathematical Functions,
``Barnes' \(G\)-function (double gamma function),'' \S5.17.
\url{https://dlmf.nist.gov/5.17}.

\bibitem{DLMF25}
NIST Digital Library of Mathematical Functions,
``Hurwitz zeta function,'' \S25.11, especially the large-parameter expansions
and their derivatives.
\url{https://dlmf.nist.gov/25.11}.

\bibitem{Edgar2010}
G.~A. Edgar,
``Transseries for beginners,''
\emph{Real Analysis Exchange} \textbf{35} (2009/2010), 253--310.

\bibitem{Ecalle1981}
J.~\'Ecalle,
\emph{Les fonctions r\'esurgentes}, Vols.~I--III,
Publications Math\'ematiques d'Orsay, 1981--1985.

\bibitem{FlajoletSedgewick2009}
P.~Flajolet and R.~Sedgewick,
\emph{Analytic Combinatorics},
Cambridge University Press, 2009.

\bibitem{Hardy1910}
G.~H. Hardy,
\emph{Orders of Infinity},
Cambridge Tracts in Mathematics and Mathematical Physics, No.~12,
Cambridge University Press, 1910.

\bibitem{Henrici1974}
P.~Henrici,
\emph{Applied and Computational Complex Analysis}, Vol.~1,
Wiley, 1974.

\bibitem{MezoBaricz2017}
I.~Mez\H{o} and \'A.~Baricz,
``On the generalization of the Lambert \(W\) function,''
\emph{Transactions of the American Mathematical Society}
\textbf{369} (2017), 7917--7934.

\bibitem{Nemes2014}
G.~Nemes,
``Error bounds and exponential improvement for the asymptotic expansion of the
Barnes \(G\)-function,''
\emph{Proceedings of the Royal Society A} \textbf{470} (2014), 20140534.
\url{https://doi.org/10.1098/rspa.2014.0534}.

\bibitem{Nemes2017}
G.~Nemes,
``Error bounds for the asymptotic expansion of the Hurwitz zeta function,''
\emph{Proceedings of the Royal Society A} \textbf{473} (2017), 20170363.
\url{https://doi.org/10.1098/rspa.2017.0363}.

\bibitem{Olver1997}
F.~W.~J. Olver,
\emph{Asymptotics and Special Functions},
A K Peters, 1997 reprint of the 1974 edition.

\bibitem{ReshetnikovPL2026}
V.~Reshetnikov,
``Polynomial--Logarithmic Transseries: Algebra, Composition, Reversion, and the
Lambert \(W\) Archetype,''
ProveIt research draft, 2026.
\url{https://github.com/VladimirReshetnikov/ProveIt/tree/main/Analysis/FabiusFunction/docs/semi-formalized-research-frontiers/drafts/series-and-transseries}.

\bibitem{ReshetnikovLambert2026}
V.~Reshetnikov,
``Asymptotic Reversion of \(x+W(x)\): A General Calculus for Logarithmic
Transseries,''
ProveIt research draft, 2026.
\url{https://github.com/VladimirReshetnikov/ProveIt/tree/main/Analysis/FabiusFunction/docs/semi-formalized-research-frontiers/drafts/series-and-transseries/lambert-inverse-transseries}.

\bibitem{vanDerHoeven2006}
J.~van der Hoeven,
\emph{Transseries and Real Differential Algebra},
Lecture Notes in Mathematics, Vol.~1888, Springer, 2006.

\bibitem{Vardi1988}
I.~Vardi,
``Determinants of Laplacians and multiple gamma functions,''
\emph{SIAM Journal on Mathematical Analysis} \textbf{19} (1988), 493--507.
\url{https://doi.org/10.1137/0519035}.

\end{thebibliography}

\end{document}
